%% ELAPSE: Entropy-Linked Autonomous Propagation and Self-Erasure
%% Target journal: Applied Mathematical Modelling (AMM), Elsevier
%% Class: elsarticle, two-column not required for AMM
%%
%% Compile: pdflatex ELAPSE_AMM.tex
%%          bibtex   ELAPSE_AMM
%%          pdflatex ELAPSE_AMM.tex  (x2)
%%
\documentclass[review,12pt]{elsarticle}

\usepackage{amsmath,amssymb,amsfonts}
\usepackage{amsthm}
\usepackage{graphicx}
\usepackage{booktabs}
\usepackage{hyperref}
\usepackage{xcolor}
\usepackage{url}
\usepackage{natbib}
\usepackage{multirow}
\usepackage{caption}
\usepackage{subcaption}
\usepackage{pifont}

%% Custom commands
\newcommand{\R}{\mathbb{R}}
\newcommand{\E}{\mathbb{E}}
\newcommand{\ind}{\mathbf{1}}       % indicator function
\newcommand{\bfx}{\mathbf{x}}
\newcommand{\bfw}{\mathbf{w}}
\newcommand{\bfv}{\mathbf{v}}
\newcommand{\IEE}{\mathrm{IEE}}
\newcommand{\Hc}{H_c}
\newcommand{\Hmax}{H_{\max}}
\newcommand{\lambdat}{\lambda_2}

\newtheorem{theorem}{Theorem}
\newtheorem{lemma}{Lemma}
\newtheorem{observation}{Observation}
\newtheorem{proposition}{Proposition}
\newtheorem{conjecture}{Empirical Conjecture}
\newtheorem{corollary}{Corollary}
\newtheorem{remark}{Remark}
\newtheorem{definition}{Definition}

\graphicspath{{../../output/figures/}}

\journal{Applied Mathematical Modelling}

\begin{document}

\begin{frontmatter}

\title{ELAPSE: Entropy-Linked Autonomous Propagation and Self-Erasure ---
       A Multi-Domain Ensemble Framework for Active Data Containment
       in Decentralised Networks}

\author[rv]{Y.~Sailaja}
\ead{sailaja@rvce.edu.in}

\author[rv]{Chhavi Pareek}
\author[rv]{Hitarth Mehra}
\author[rv]{Harsh Patel}

\address[rv]{Department of Mathematics / Department of Computer Science \&
             Engineering, RV College of Engineering, Bangalore 560059, India}

\begin{abstract}
Restricting information flow across stochastic, decentralised peer-to-peer
networks is a critical open challenge at the intersection of network science,
stochastic analysis, and data privacy engineering. Classical cryptographic
containment protocols~\cite{geambasu2009vanish,tang2010fade} trigger
self-erasure on chronological timers, rendering them blind to the actual
topological spread of a data payload. As a result, scale-free networks ---
in which viral diffusion can compromise a majority of nodes in a small
fraction of any fixed timeout window --- systematically defeat
time-based deletion guarantees.

We introduce \textbf{ELAPSE} (Entropy-Linked Autonomous Propagation and
Self-Erasure), a framework that replaces the chronological clock with the
Shannon entropy of the data's spatial dispersion as the mortality trigger.
Five constituent mechanisms drawn from mathematical physics, epidemiology,
stochastic finance, systems biology, and sociological cascade theory each
cast a continuous vote $v_k(t) \in [0,1]$ measuring how strongly conditions
warrant deletion. These votes are fused by an entropy-regularised Nelder-Mead
ensemble that minimises the Time-Integrated Entropy Exposure (IEE), a
functional tracking both spatial privacy loss and surviving data mass.

We make five primary contributions.
\emph{First}, we prove a suite of theoretical results: a well-posedness theorem
(Yamada-–Watanabe); a closed-form spectral IEE bound; a topology-refined
refinement connecting the spectral trigger time to the algebraic connectivity
$\lambda_2$; a residual-mass privacy guarantee showing that IEE control implies
exponential data-recovery bounds for bounded adversaries; and exact convexity
of the linearised IEE objective via a Laplace-transform argument.
\emph{Second}, we provide statistically rigorous results: all IEE values
are reported as mean $\pm$ 95\% confidence intervals over $N_{\mathrm{test}}=20$
out-of-sample test seeds per (topology, network-size, mechanism) combination.
\emph{Third}, we validate on real-world P2P networks from the Stanford SNAP
repository~\cite{snapnets} (Gnutella08 and Gnutella31 subgraphs), confirming
that synthetic findings transfer to empirical topologies.
\emph{Fourth}, we introduce adversarial and distributed-estimation extensions:
an attacker model in which a fraction $f$ of nodes throttle propagation to
evade entropy-triggered deletion, together with an adaptive-$\Hc$
countermeasure; and a gossip-protocol approximation to the global entropy
observable that requires only $k$-hop neighbourhood information.
\emph{Fifth}, we characterise the minimum gossip hop count $k^*$ as a
function of network algebraic connectivity, providing practitioners with
topology-specific deployment guidance.

Numerical experiments across Erd\H{o}s-R\'enyi, Barab\'asi-Albert, and
Watts-Strogatz topologies ($n \in \{50, 100, 200, 500\}$) confirm that the
ELAPSE ensemble achieves IEE at least $9\times$ lower than the uncontrolled
baseline at $n = 500$.
Under ER-only weight training (primary evaluation at $n=500$), Social Cascade
(M5) achieves lower per-topology IEE on all three topologies, with the ensemble
trailing M5 by $3\%$ on ER and ${\approx}33\%$ on BA and WS due to
weight-transfer failure; topology-matched retraining closes the gap to within
$5$--$8\%$ on BA/WS and yields M6 outperforming M5 by $5.9\%$ on ER ($n=100$),
achieving the lowest topology-averaged IEE across all mechanisms tested.
The ensemble's principal advantages over any single mechanism are
\emph{robustness to threshold miscalibration} and \emph{adversarial resilience}:
ELAPSE strictly dominates the bang-bang policy (M8) when $\Hc$ exceeds the
observed peak entropy --- the practically common case of conservative threshold
setting --- and limits IEE degradation under hub-throttling attacks to
${\leq}6.6\%$, while the same attack drives M8 to complete deletion failure.
\end{abstract}

\begin{keyword}
Decentralised networks \sep Information diffusion \sep Shannon entropy \sep
Self-erasure \sep Active data containment \sep Stochastic differential equations
\sep Ensemble optimisation \sep Right to erasure \sep GDPR
\end{keyword}

\end{frontmatter}

%% ─── 1. INTRODUCTION ───────────────────────────────────────────────────────
\section{Introduction}
\label{sec:intro}

\subsection{Motivation and Problem Statement}

The proliferation of sensitive data across decentralised peer-to-peer (P2P)
networks constitutes a qualitatively different privacy threat than the storage
of data in centralised repositories. Once a data payload escapes its point of
origin --- whether via legitimate forwarding, an insider disclosure, or
endpoint compromise --- the originating system loses all physical control over
its further replication. The challenge is acute in the context of the European
Union's \emph{right to erasure} (GDPR Article 17~\cite{gdpr2016}), which mandates
that data subjects can demand removal of their personal data; a mandate that is
straightforwardly unenforceable once a payload has diffused across thousands
of pseudonymous P2P nodes.

The scale of the problem is growing. IBM's 2024 Cost of a Data Breach
Report~\cite{ibm2024breach} documented a global average breach cost of
USD~4.88~million, a 10\% year-on-year increase, with P2P exfiltration cited
among the leading attack vectors. The 2024 Verizon Data Breach Investigations
Report~\cite{verizon2024dbir} further notes that data discovered in unauthorised
copies on external networks has a median ``time to identification'' exceeding
280~days, during which the privacy exposure accumulates continuously.

\subsection{Limitations of Chronological Containment}

The existing active-containment literature --- most influentially the
Vanish~\cite{geambasu2009vanish} and FADE~\cite{tang2010fade} systems --- relies
on a shared architectural assumption: a \emph{chronological decay timer} governs
data lifetime. Vanish overwrites Distributed Hash Table (DHT) shards on a fixed
schedule (typically 8~hours), so that the decryption key is irretrievably lost
after a wall-clock timeout. FADE extends this with file-assured deletion via
policy-based key management.

The fundamental limitation of this paradigm is its independence from the actual
network topology. On a scale-free (Barab\'asi-Albert) graph, the hub-driven
spreading mechanism can saturate a majority of nodes within a small fraction of
any fixed timeout~\cite{barabasi1999emergence,pastor2015epidemic}. A data payload
that has already reached $10^3$ nodes before the 8-hour timer fires may be
practically irrecoverable regardless of key deletion. Conversely, on a
sparse small-world (Watts-Strogatz) graph, the same timer may fire when the
data has barely left the source neighbourhood, over-triggering deletion.

\subsection{The Entropy-Clock Approach}

We propose that data lifetime should be a \emph{dynamic function of the
network state}. Specifically, ELAPSE uses the Shannon entropy of the data's
spatial distribution as its deletion clock:
\begin{equation}
  H(t) = -\sum_{i=1}^{n} p_i(t) \ln p_i(t), \quad
  p_i(t) = \frac{x_i(t)}{\sum_j x_j(t)},
  \label{eq:entropy}
\end{equation}
where $x_i(t) \geq 0$ is the data concentration at node $i$ at time $t$.
When $H(t)$ is low, data is concentrated on few nodes (low privacy exposure);
when $H(t)$ approaches $\Hmax = \ln n$, data has spread uniformly (maximal
exposure). Self-erasure is triggered when $H(t)$ crosses a critical
threshold $\Hc$, replacing the chronological clock with a
\emph{spatial-topological clock}.

\subsection{Contributions}

This paper makes the following contributions:
\begin{enumerate}
  \item \textbf{Framework}: We introduce ELAPSE, fusing five mechanism-specific
        deletion votes via an entropy-regularised Nelder-Mead ensemble optimising
        IEE (Section~\ref{sec:mechanisms}--\ref{sec:ensemble}).
  \item \textbf{Rigorous theory} (Section~\ref{sec:theory}):
        Five results: (i) pathwise-unique non-negative strong solution
        (Theorem~\ref{thm:wellposed}, via Yamada-–Watanabe);
        (ii) a closed-form spectral IEE upper bound in terms of $n$, $M_0$,
        and deletion intensity, valid for the \emph{deterministic limiting
        system} in the post-trigger phase (Theorem~\ref{thm:iee_bound};
        see Remark~\ref{rem:det_bound} for the stochastic scope);
        (iii) a \emph{topology-refined} IEE bound incorporating the
        algebraic connectivity $\lambda_2$, giving the spectral trigger
        time $\tau(\lambda_2)$ explicitly (Corollary~\ref{cor:lambda2});
        (iv) a residual-mass privacy guarantee connecting IEE to an
        exponential data-recovery bound for uniformly-random bounded adversaries
        (Proposition~\ref{prop:privacy}; hub-targeted adversary bound is open);
        (v) \emph{exact} convexity of the \emph{linearised} IEE objective
        via a Laplace-transform argument (Theorem~\ref{thm:lin_convex},
        Corollary~\ref{cor:unique}); the full nonlinear SDE has a 14.1\%
        Jensen-violation rate, quantifying the linearisation gap.
  \item \textbf{Statistical rigour}: All IEE values are reported as
        mean $\pm$ 95\% CI over $N_{\mathrm{test}} = 20$ out-of-sample seeds,
        three topologies, and four network sizes (Section~\ref{sec:results}).
  \item \textbf{Real-world validation}: The full ensemble is evaluated on
        hub-seeded BFS subgraphs of the Stanford SNAP Gnutella P2P
        datasets~\cite{snapnets} (Section~\ref{sec:snap}).
        This primarily validates BA-regime performance on an independent
        empirical dataset; ER- and WS-regime validation on real networks
        remains an open empirical challenge (see Section~\ref{sec:snap_topology}).
  \item \textbf{Adversarial analysis}: We model an attacker who throttles
        propagation and propose an adaptive-$\Hc$ countermeasure
        (Section~\ref{sec:adversarial}).
  \item \textbf{Distributed estimation}: We propose a gossip-protocol
        approximation to the global entropy observable and quantify its
        trigger-accuracy (available as supplementary material).
\end{enumerate}

\subsection{Paper Organisation}

Section~\ref{sec:related} reviews the literature and positions ELAPSE.
Section~\ref{sec:mechanisms} defines the five core mechanisms.
Section~\ref{sec:ensemble} presents the ensemble construction and five
theoretical results (Theorems~\ref{thm:wellposed}--\ref{thm:lin_convex},
Corollaries~\ref{cor:lambda2}--\ref{cor:unique},
Propositions~\ref{prop:privacy}--\ref{prop:vs_timer}).
Section~\ref{sec:results} provides the main statistical evaluation.
Section~\ref{sec:snap} reports real-world validation.
Section~\ref{sec:adversarial} analyses adversarial evasion.
Section~\ref{sec:conclusions} concludes.
Appendix~\ref{sec:m3} provides a mechanistic analysis and fix for
M3's underperformance.

%% ─── 2. RELATED WORK ───────────────────────────────────────────────────────
\section{Related Work}
\label{sec:related}

\subsection{Information Diffusion and Entropy on Graphs}

The quantitative study of information propagation on complex networks has a
rich lineage. Foundational epidemiological models for spreading
processes~\cite{pastor2015epidemic} established that network topology --- and
especially heterogeneous degree distributions --- fundamentally alters
threshold phenomena. The particular role of Shannon entropy as a topological
indicator gained prominence through Zhao et al.~\cite{zhao2024multi}, who
demonstrated that entropy-spreading metrics provide accurate macroscopic proxies
for information volatility across heterogeneous degree distributions, establishing
$H(t)$ as a reliable outbreak-detection signal.

Parallel advances in continuous graph diffusion, originating in spectral graph
signal processing~\cite{shuman2013emerging} and extended through neural
PDE-on-graph formulations (GRAND~\cite{chamberlain2021grand};
Eliasof et al.~\cite{eliasof2024graph}; Fan et al.~\cite{fan2024graph}), confirm
that partial differential equations on graph topologies accurately capture
non-linear propagation boundaries. Universal resilience patterns~\cite{gao2016universal}
further characterise scale-stress behaviour under perturbation.

\subsection{Static and Active Data Containment}

Early cryptographic containment schemes were pioneered by
Vanish~\cite{geambasu2009vanish}, which exploited DHT key churn to implement
passive key decay. Wolchok et al.~\cite{wolchok2010defeating} demonstrated
Sybil-attack vulnerability. FADE~\cite{tang2010fade} and the
Ephemerizer~\cite{perlman2005ephemerizer} strengthened storage layers but
retained chronological timers. Rondeau and Temple~\cite{rondeau2021secure}
survey secure deletion approaches and confirm that time-blind mechanisms are
structurally insufficient for dynamic network environments. Truex et al.~\cite{truex2019hybrid} demonstrate a hybrid local/global
approach to privacy-preserving federated learning in distributed settings,
showing that no single privacy mechanism dominates across heterogeneous
network topologies — a finding consistent with our ensemble design.

\subsection{Machine Unlearning and Algorithmic Forgetting}

The machine-learning community has recently formalised the complementary
problem of ``algorithmic forgetting'': ensuring a trained model provides no
residual information about a removed training sample. Bourtoule
et al.~\cite{bourtoule2021machine} established the SISA training framework,
Nguyen et al.~\cite{nguyen2025survey} and Xu et al.~\cite{xu2023machine}
provide comprehensive surveys of unlearning algorithms, highlighting the
need for topological removal verification. Oesterling
et al.~\cite{oesterling2024fair} incorporate fairness constraints into active
data removal, which motivates our entropy-regularised ensemble approach.

\subsection{Differential Privacy and Information-Theoretic Bounds}

Dwork and Roth's foundational treatise~\cite{dwork2014algorithmic} establishes
differential privacy as the canonical framework for information-theoretic
privacy guarantees. While differential privacy bounds \emph{inference
leakage} from a fixed dataset, ELAPSE addresses a complementary problem:
minimising the \emph{spatial exposure} of a diffusing payload before deletion
can take effect. These two regimes are formally distinct but share the use of
entropy-based metrics as core sensitivity measures.

\subsection{Control of Diffusion Processes on Graphs}

ELAPSE is related to, but distinct from, the emerging literature on
\emph{controlled diffusion on networks}.
Mean-field game (MFG) frameworks~\cite{carmona2018probabilistic,achdou2020mean}
study populations of McKean-–Vlasov-type agents interacting through their
empirical distribution, deriving feedback control laws from coupled
Hamilton-–Jacobi-–Bellman and Fokker-–Planck equations.
In MFG language, ELAPSE can be interpreted as a \emph{social planner}
problem: we seek a mortality schedule $\mu\Delta(t;\bfw)$ that minimises
a population-level cost (IEE) subject to the coupled McKean-–Vlasov SDE
\eqref{eq:sde}.  However, three features distinguish ELAPSE from classical
MFG: (i)~the graph topology enters via the Laplacian $L$ rather than a
mean-field coupling term; (ii)~the control $\Delta(t;\bfw)$ is an
\emph{entropy-dependent} functional of the empirical distribution, not an
agent-specific feedback; and (iii)~we optimise weights over a finite
ensemble of pre-specified mechanisms rather than solving a PDE.
Online feedback control on network flow
manifolds~\cite{bianchin2021online} addresses a related setting (closed-loop
optimisation on constrained graph dynamics) but focuses on power systems
rather than information-theoretic objectives.
The present work contributes to this space by providing the first
entropy-functional IEE bound (Theorem~\ref{thm:iee_bound}) and spectral
trigger time (Corollary~\ref{cor:lambda2}) for controlled diffusion on
heterogeneous P2P graphs.

\paragraph{Comparison baselines: scope and limitations.}
The present work evaluates ELAPSE against the chronological timer (M7)
and the greedy bang-bang policy (M8).  We do not compare against
model-predictive control (MPC) or reinforcement-learning (RL) deletion
policies, both of which are natural closed-loop alternatives.  MPC would
optimise deletion intensity over a rolling horizon using the SDE model as
a predictor; RL would train a policy on the reward $-H(t)M(t)$ without
an explicit model.  These comparisons are deferred to future work
(Section~\ref{sec:conclusions}) for two reasons: (i) MPC requires solving
the SDE forward at each time step, which is computationally expensive
at the network sizes studied here; and (ii) RL baseline training would
require substantially larger computational resources to achieve stable
convergence across all three topology classes and four network sizes.
The bang-bang policy (M8), which is the optimal myopic controller for
the linearised SDE, serves as the strongest tractable baseline.

\subsection{Positioning of ELAPSE}

Table~\ref{tab:positioning} summarises how ELAPSE differs from the three
dominant paradigms for data-lifetime management.

\begin{table}[htbp]
  \caption{Comparison of data-deletion paradigms.  ELAPSE is the only
           approach that is simultaneously topology-aware, operates at
           network level, minimises a formal IEE objective, and provides
           distributed estimation.}
  \label{tab:positioning}
  \centering
  \small
  \renewcommand{\arraystretch}{1.15}
  \begin{tabular}{lcccc}
    \toprule
    Paradigm & Topology- & Network- & Formal & Distributed \\
             & aware     & level    & IEE obj. & estimation \\
    \midrule
    Cryptographic (Vanish, FADE)~\cite{geambasu2009vanish,tang2010fade}
              & \ding{55} & \ding{55} & \ding{55} & \ding{55} \\
    Chron.\ timer (M7, \S\ref{sec:m7_comparison})
              & \ding{55} & \ding{55} & \ding{55} & \ding{55} \\
    Machine unlearning~\cite{bourtoule2021machine,xu2023machine}
              & \ding{55} & \ding{55} & \ding{55} & \ding{55} \\
    Differential privacy~\cite{dwork2014algorithmic}
              & \ding{55} & \ding{55} & \ding{55} & Partial \\
    \textbf{ELAPSE (ours)}
              & \ding{51} & \ding{51} & \ding{51} & \ding{51} \\
    \bottomrule
  \end{tabular}
\end{table}

ELAPSE occupies a unique position in the landscape: it is the first framework
that (i) replaces the chronological erasure clock with a spatial entropy
trigger, (ii) constructs a regularised ensemble across five physical domains,
(iii) provides rigorous theoretical guarantees (Theorems~\ref{thm:wellposed}--\ref{thm:lin_convex})
backed by statistical evaluation, and (iv) characterises the minimum
gossip hop count $k^*$ as a function of network algebraic connectivity
(Proposition~\ref{prop:gossip_lambda}).
Prior ensemble approaches for network spreading
(e.g., competing epidemics~\cite{karrer2011competing}; social
cascades~\cite{guilbeault2018complex}) do not address active containment.

%% ─── 3. FIVE CORE MECHANISMS ───────────────────────────────────────────────
\section{The Five Core Mechanisms}
\label{sec:mechanisms}

\subsection{Network and Notation}

Let $G = (V, E)$ be an undirected connected graph with $|V| = n$ nodes. The
graph Laplacian $L = D - A$ (where $D$ is the degree matrix and $A$ the
adjacency matrix) and the Fiedler value $\lambdat$ (second-smallest eigenvalue
of $L$, the algebraic connectivity~\cite{chung1997spectral}) play central roles.
The data concentration at node $i$ at time $t$ is $x_i(t) \geq 0$, with units
of \emph{data mass per node}.  Initial conditions are $x_i(0) = u_i$ where
$u_i \sim \mathrm{Uniform}(0.5, 1.0)$ for a randomly chosen 10\% of nodes
(the initial data sources) and $x_i(0) = 0$ otherwise; total initial mass
$M(0) = \sum_i x_i(0) \in (0, n/10]$.  All IEE values therefore have units
of [nats $\times$ data-mass $\times$ time], and scale with both $n$ and $M(0)$.
Spatial privacy exposure is measured by $H(t)$ as defined in Eq.~\eqref{eq:entropy},
with $\Hmax = \ln n$ and a critical threshold $\Hc = 0.65 \Hmax$ used throughout.

Each of the five mechanisms defines a continuous \emph{deletion vote}
$v_k(t) \in [0,1]$, measuring urgency of erasure. The composite mortality
pressure $\Delta(t) = \sum_k w_k v_k(t)$ drives data removal.

\subsection{M1: Entropy-Gated Diffusion Model (EGDM)}

The physical baseline treats data as a field propagating against a
graph-Laplacian operator with an entropy-triggered global mortality term:
\begin{equation}
  \frac{d\bfx}{dt} = -\alpha L\bfx - \mu\,\sigma(H(t))\,\bfx
                     + \mathbf{s}(t) + dW_t,
  \label{eq:egdm}
\end{equation}
where $\alpha = 0.3$ is the diffusion coefficient, $\mu = 1.5$ the mortality
rate, $\mathbf{s}(t)$ a small external injection, and $dW_t$ a
multiplicative CIR-type Wiener noise with scale $\sigma_n = 0.015$.
The polynomial vote is:
\begin{equation}
  v_1(t) = \left[\max\!\left(0,\;\frac{H(t) - \Hc}{\Hmax - \Hc}\right)\right]^{\!\beta},
  \label{eq:v1}
\end{equation}
with sharpness $\beta = 2$. Near-zero below $\Hc$, the vote rises as a
smooth power-law as data spreads. Strong Lipschitz continuity of the
deterministic flow ensures existence and uniqueness of strong solutions
up to the erasure threshold~\cite{oksendal2003stochastic}.

\subsection{M2: Epidemic-Entropy (Compartmental SIR)}

Epidemiology models data spread as an infectious load. Nodes transition
through Susceptible-Infected-Recovered compartments, capturing local
neighbour-to-neighbour spreading independently of global diffusion:
\begin{equation}
  \frac{dI_i}{dt} = \beta_{\mathrm{sir}}\,S_i \sum_j A_{ij} I_j
                   - \gamma I_i - \mu\,\sigma(H(t))\,I_i,
  \label{eq:sir}
\end{equation}
where $\beta_{\mathrm{sir}} = 0.4$, $\gamma = 0.15$. The vote is the
current infected fraction: $v_2(t) = n^{-1}\!\sum_i I_i(t)$. This mechanism
fires rapidly on dense-core topologies but is suppressed by recovery on
sparse graphs.

\subsection{M3: Stochastic-Finance (Ornstein-Uhlenbeck First-Passage)}

Drawing from barrier-option pricing~\cite{black1973pricing,uhlenbeck1930theory},
information exposure risk is modelled as a mean-reverting stochastic process
approaching a ruin boundary. The state evolves as:
\begin{equation}
  d\bfx = \theta(\mu_{ou} - \bfx)\,dt + \sigma_{ou}\,dW_t
         - \mu\,P(\tau \leq t)\,\bfx\,dt,
  \label{eq:finance}
\end{equation}
with $\theta = 0.3$, $\mu_{ou} = 0.1$, $\sigma_{ou} = 0.04$. The vote
is the first-passage probability:
\begin{equation}
  v_3(t) = 1 - \exp\!\left(-\lambda_{\mathrm{fp}}\,t\right), \quad
  \lambda_{\mathrm{fp}} = \frac{\alpha\lambdat}{\Hmax - \Hc},
  \label{eq:v3}
\end{equation}
encoding the spectral geometry of the network through the Fiedler
value $\lambdat$. As analysed in Section~\ref{sec:m3}, the static use
of $\lambdat$ limits M3's responsiveness to fast BA hub dynamics.

\begin{remark}[M3 underperformance and recommended fix]
\label{rem:m3_fix}
The static Fiedler value $\lambdat$ in \eqref{eq:v3} causes $v_3$ to rise
monotonically at a fixed rate, which poorly tracks rapid entropy growth on
BA (hub-driven) topologies where $\lambdat$ is small.  A time-varying proxy
$\lambdat(t) \approx \lambdat(0)\cdot(H(t)/\Hmax)\cdot(M(t)/M_0)$
reduces M3 standalone IEE on BA by $6.6\%$ ($p<0.01$); the static
formulation is retained here for consistency with reported experiments,
but the time-varying variant is recommended for production deployments.
The mechanistic analysis is in Appendix~\ref{sec:m3}.
\end{remark}

\subsection{M4: Biological Cooperative Binding (Hill Function)}

In systems biology, cooperative binding transitions are modelled by the
Hill function~\cite{alon2019systems}:
\begin{equation}
  v_4(t) = \frac{H(t)^{n_h}}{H_c^{n_h} + H(t)^{n_h}},
  \label{eq:hill}
\end{equation}
with Hill coefficient $n_h = 4$. This produces a smoother, more
Lyapunov-stable asymptotic transition than the polynomial $v_1$: it is
near zero for $H \ll \Hc$, exactly $\tfrac{1}{2}$ at $H = \Hc$, and
near one for $H \gg \Hc$. The sigmoidal shape is mathematically
well-characterised from gene regulatory network analysis.

\subsection{M5: Sociological Cascade Pressure}

Sociological cascade models parameterise agents that delete based on
localised sub-graph density~\cite{watts2002simple,guilbeault2018complex}.

\begin{definition}[Indicator notation]
\label{def:indicator}
For any predicate $P$, let $\ind[P]$ denote the indicator function:
$\ind[P] = 1$ if $P$ is true, $\ind[P] = 0$ otherwise.
Specifically, for a threshold $\kappa$ and deletion threshold $\epsilon$:
$\ind[x_j < \epsilon]$ equals 1 if node $j$ has effectively deleted its
payload (concentration below $\epsilon = 0.05$), and 0 otherwise;
$\ind[\cdot > \kappa]$ equals 1 if the argument exceeds the cascade
threshold $\kappa = 0.3$.
\end{definition}

The cascade vote combines global and local deletion pressure:
\begin{equation}
  v_5(t) = \tfrac{1}{2}\sigma\!\bigl(H(t)\bigr)
           + \tfrac{1}{2n}\sum_{i=1}^{n}
  \ind\!\!\left[
    \frac{\sum_{j \in N(i)}\ind[x_j < \epsilon]}{\deg(i)}
    > \kappa
  \right],
  \label{eq:social}
\end{equation}
where $N(i)$ is the neighbourhood of node $i$ and $\deg(i)$ its degree.
A node responds to deletion when the global entropy exceeds $\Hc$
\emph{or} when more than $\kappa$ of its neighbours have already deleted.

%% ─── 4. ENSEMBLE CONSTRUCTION ─────────────────────────────────────────────
\section{The ELAPSE Ensemble}
\label{sec:ensemble}

\subsection{Ensemble Architecture}

The M6 ELAPSE engine fuses the five votes into a terminal deletion pressure:
\begin{equation}
  \Delta(t) = \bfw^\top \bfv(t), \quad
  \bfw \in \mathcal{W} = \!\left\{\bfw \in \R^5_{\geq 0} : \|\bfw\|_1 = 1\right\},
  \label{eq:ensemble}
\end{equation}
where $\bfv(t) = [v_1(t),\ldots,v_5(t)]^\top$.
The Time-Integrated Entropy Exposure (IEE) is:
\begin{equation}
  \IEE(\bfw) = \int_0^T H\!\bigl(\mathbf{p}(t;\bfw)\bigr)\cdot M(t;\bfw)\,dt,
  \label{eq:iee}
\end{equation}
where $M(t;\bfw) = \sum_i x_i(t)$ is the surviving data mass. A lower
IEE indicates better privacy protection.

\begin{remark}[Metric sensitivity]
\label{rem:metric}
IEE weights entropy and mass equally in the product $H(t)\cdot M(t)$.
A scenario with high $H$ but low $M$ (data spread thin, mostly deleted)
differs in practical risk from low $H$ with high $M$ (data concentrated
but still present).  We note three robustness properties.
\emph{(i) Entropy-only weighting} $(\int_0^T H\,dt)$: all mechanisms
reduce entropy monotonically; mechanism rankings are preserved.
\emph{(ii) Mass-only weighting} $(\int_0^T M\,dt)$: ELAPSE still achieves
$\geq 9\times$ reduction over M0 at $n=500$, and the ER timing advantage
over M7 persists since early deletion directly reduces $M$.
\emph{(iii) Power-law IEE} $(\int H^\alpha M^{1-\alpha}\,dt)$: by
continuity in $\alpha$, mechanism rankings are preserved for
$\alpha \in [0.3, 0.7]$ near the symmetric product $(\alpha = 0.5)$.
A systematic metric sensitivity analysis with Sobol indices across $\alpha$
is identified as future work.
\end{remark}

\subsection{Linear Vote Fusion: Derivation and Alternatives}
\label{sec:fusion_justification}

The linear combination $\Delta(t) = \bfw^\top\bfv(t)$ is the natural
parametrisation of additive-mortality control in the ELAPSE SDE.

\begin{proposition}[Linear fusion is the canonical implementation of additive-mortality control]
\label{prop:fusion}
In the ELAPSE SDE, mortality enters as $-\mu\Delta(t;\bfw)\,x_i$ with
$\Delta(t;\bfw) = \bfw^\top\bfv(t)$, so that mechanism $k$ contributes
deletion rate $\mu w_k v_k(t)$ independently of other mechanisms.
This additive decomposition is exact by construction: there are no
cross-mechanism interaction terms in the SDE.  Consequently,
$w_k$ directly and exclusively scales mechanism $k$'s mortality contribution,
making the weight space $\mathcal{W}$ a minimal, identifiable parametrisation
of the ensemble control authority.  Any nonlinear reparametrisation
$\tilde{w}_k = g(w_k)$ would couple mechanisms through $g$, conflating
their individual contributions and preventing direct attribution of IEE
changes to individual mechanism weights.
\end{proposition}

\begin{remark}[Scope of the proposition]
Proposition~\ref{prop:fusion} justifies linear fusion as \emph{structurally
correct} for additive SDE mortality, not as universally optimal among all
conceivable fusion rules.  Whether nonlinear fusion (e.g., product-of-experts
or max-vote) achieves lower IEE for some SDE with interaction terms is a
separate question; within the ELAPSE architecture, additivity is not an
assumption but a consequence of the SDE design.
\end{remark}

\paragraph{Comparison with nonlinear fusion rules.}
Two natural alternatives are the \emph{product-of-experts}
$\Delta_{\rm PoE}(t) = \prod_k v_k(t)^{w_k}$
(geometric mean) and the \emph{maximum-vote}
$\Delta_{\rm max}(t) = \max_k (w_k v_k(t))$.
Product-of-experts requires all mechanisms to simultaneously vote for deletion
(logical AND), which is over-conservative and delays $\Delta$ even when four
of five mechanisms signal high entropy.  Maximum-vote uses only the single
most-confident mechanism, ignoring the diversity of evidence.  The linear
combination provides a smooth interpolation: it fires when the
\emph{weighted majority} of mechanisms vote positively.
Quantitative IEE comparison of these fusion classes is identified as future work,
as it requires independent ensemble optimisation for each class.

\subsection{Entropy-Regularised Optimisation}

Unconstrained Nelder-Mead on $\mathcal{W}$ collapses to a one-hot vertex
$\mathbf{e}_k$ (choosing the single best mechanism), sacrificing
generalisation across topology types. We add an entropy regulariser:
\begin{equation}
  \min_{\bfw \in \mathcal{W}}\; J(\bfw)
    = \IEE(\bfw) - \lambda_{\mathrm{reg}}\,\mathcal{H}(\bfw),
  \label{eq:objective}
\end{equation}
where $\mathcal{H}(\bfw) = -\sum_k w_k \ln w_k$ is the weight-layer
entropy. Subtracting this term penalises low-entropy (collapsed) weight
distributions while allowing IEE minimisation to proceed. We set
$\lambda_{\mathrm{reg}} = 15$ throughout, selected at the Pareto elbow
between IEE and weight diversity on BA topology at $n=200$
(see Table~\ref{tab:reg}); cross-topology validation of this hyperparameter
is identified as future work.

\subsection{Theoretical Analysis}
\label{sec:theory}

We establish three main results: well-posedness of the SDE system, a
closed-form spectral upper bound on IEE, and rigorous convexity of the
linearised IEE objective.

\begin{theorem}[Well-posedness of ELAPSE]
\label{thm:wellposed}
Let $G$ be connected and $\bfx(0) \geq 0$ componentwise.
The ELAPSE SDE system \eqref{eq:egdm}--\eqref{eq:ensemble} has a
pathwise-unique, non-negative strong solution $\bfx(t) \geq \mathbf{0}$
a.s.\ on $[0,T]$ for any finite $T > 0$.
\end{theorem}

\begin{proof}
We establish the four required properties in sequence.

\textbf{(i) Local existence via locally Lipschitz drift.}
Each vote $v_k(t;\bfx) \in [0,1]$ is a continuous function of $H(\bfx)$, and
$H$ is $C^1$ on the open set $\Omega_c = \{\bfx : M(\bfx) \geq c > 0\}$.
On any compact $K \subset \Omega_c$, the drift
$b_i(\bfx,t) = -\alpha(L\bfx)_i - \mu\Delta(t;\bfw)x_i + s_i(t)$
is therefore Lipschitz in $\bfx$ with constant $L_K = \alpha\|L\|_2 + \mu + \|\nabla_\bfx H\|_{C(K)}$.
Local existence and uniqueness on $[0,\tau_K)$ (up to the first exit from $K$)
follows from the Picard--Lindelöf theorem for SDEs~\cite{oksendal2003stochastic}.

\textbf{(ii) Global existence via Lyapunov bound.}
Let $V(\bfx) = M(\bfx) = \sum_i x_i$.  Itô's formula gives:
\[
  dM = \Bigl(-\mu\Delta(t;\bfw)M + \|{\bf s}\|_1\Bigr)dt
       + \sigma_n\sum_i \sqrt{x_i}\,dW_i,
\]
using $\mathbf{1}^\top L = \mathbf{0}$.  Taking expectations:
$d\mathbb{E}[M]/dt \leq \|{\bf s}\|_1$,
so $\mathbb{E}[M(t)] \leq M_0 + \|{\bf s}\|_1 T < \infty$ for all $t \in [0,T]$.
By Markov's inequality the solution does not exit $\Omega_c$ in finite expected
time, ruling out blow-up and allowing extension to $[0,T]$.

\textbf{(iii) Non-negativity.}
Fix any component $i$.  When $x_i = 0$ the diffusion coefficient
$\sigma_n\sqrt{x_i}$ vanishes exactly, and the drift reduces to
$b_i\!\mid_{x_i=0} = \alpha\sum_j A_{ij}x_j + s_i \geq 0$ (inward
Laplacian flux plus source injection).  The component $x_i$ therefore
satisfies the CIR-type SDE
$dx_i \geq (\alpha\sum_j A_{ij}x_j + s_i)\,dt + \sigma_n\sqrt{x_i}\,dW_i$
at the boundary.  By the comparison theorem for CIR processes
(Feller~\cite{feller1951two}; see also Theorem~IX.3.7 in~\cite{revuzyor1999}),
$Y \equiv 0$ is a strict lower barrier: $x_i(t) \geq 0$ a.s.\ on $[0,T]$.

\textbf{(iv) Pathwise uniqueness via Yamada--Watanabe.}
The diffusion coefficient $\sigma(u) = \sigma_n\sqrt{u}$ satisfies
$|\sigma(u) - \sigma(v)| = \sigma_n|\sqrt{u}-\sqrt{v}|$.
For $u,v \geq 0$: $(\sqrt{u}-\sqrt{v})^2(\sqrt{u}+\sqrt{v})^2 = (u-v)^2$,
so $|\sqrt{u}-\sqrt{v}| = |u-v|/(\sqrt{u}+\sqrt{v}) \leq \sqrt{|u-v|}$,
giving $|\sigma(u)-\sigma(v)| \leq \sigma_n\sqrt{|u-v|}$.
With $h(r)=\sigma_n\sqrt{r}$: $\int_0^1 h(r)^{-2}\,dr = \sigma_n^{-2}\int_0^1 r^{-1}\,dr = +\infty$,
verifying the Yamada--Watanabe Hölder-$\tfrac{1}{2}$ condition
(Theorem~IX.3.5 in~\cite{revuzyor1999}).
Combined with (i)--(iii), the Yamada--Watanabe theorem
yields a pathwise-unique strong solution.
\end{proof}

\begin{lemma}[Positivity of deletion intensity]
\label{lem:delta_pos}
Let $\bfw^* \in \mathrm{int}(\mathcal{W})$ be the ensemble weight vector
guaranteed by Corollary~\ref{cor:unique}.  Define
$\tau^* = \tau(\lambdat)$ from \eqref{eq:tau_lambda}.  Then:
\begin{equation}
  \delta_{\min} \;:=\; \inf_{t \in [\tau^*, T]} \Delta(t;\bfw^*) \;>\; 0.
  \label{eq:delta_min}
\end{equation}
\end{lemma}

\begin{proof}
Corollary~\ref{cor:lambda2} gives $H(t) \geq \Hc$ for all $t \geq \tau^*$.
Since $H(t) > \Hc$ on a set of positive measure (strict inequality holds
except at the measure-zero boundary), the M1 vote satisfies
$v_1(t) = [(H(t)-\Hc)/(\Hmax-\Hc)]^\beta > 0$ on $(\tau^*,T]$.
As $\bfw^* \in \mathrm{int}(\mathcal{W})$ we have $w_1^* > 0$, so
$\Delta(t;\bfw^*) \geq w_1^* v_1(t) > 0$.
Since $v_1$ is continuous and strictly positive on the compact set $[\tau^*,T]$,
its infimum is strictly positive; hence $\delta_{\min} > 0$.
\end{proof}

\begin{theorem}[Spectral IEE bound]
\label{thm:iee_bound}
Let $\delta_{\min} > 0$ be as in Lemma~\ref{lem:delta_pos}
(positivity guaranteed by Corollary~\ref{cor:unique}).
Let $\tau^* = \tau(\lambdat)$ from \eqref{eq:tau_lambda} be the spectral
trigger time.  The IEE of the controlled system satisfies:
\medskip

\noindent\emph{(a) Pre-trigger phase $[0, \tau^*]$:}
$\displaystyle\int_0^{\tau^*} H(t)M(t)\,dt \leq \ln(n)\,M_0\,\tau^*$.

\noindent\emph{(b) Post-trigger phase $[\tau^*, T]$:} For the deterministic
system ($\sigma_n=0$, $\mathbf{s}\equiv\mathbf{0}$):
\begin{equation}
  \int_{\tau^*}^T H(t)M(t)\,dt \;\leq\;
  \ln(n)\cdot M_0 \cdot
  \frac{1 - e^{-\mu\delta_{\min}(T-\tau^*)}}{\mu\delta_{\min}}.
  \label{eq:iee_bound}
\end{equation}
Combining (a) and (b):
\begin{equation}
  \IEE(\bfw) \;\leq\; \ln(n)\,M_0\!\left[
    \tau^* + \frac{1-e^{-\mu\delta_{\min}(T-\tau^*)}}{\mu\delta_{\min}}
  \right].
  \label{eq:iee_total}
\end{equation}
\end{theorem}

\begin{proof}
\emph{Part (a):} Since $H(t) \leq \ln n$ and $M(t) \leq M_0$ trivially:
$\int_0^{\tau^*} H M\,dt \leq \ln(n)\,M_0\,\tau^*$.
This bound holds for both the stochastic and deterministic systems.

\emph{Part (b):} On $[\tau^*,T]$, Lemma~\ref{lem:delta_pos} gives
$\Delta(t;\bfw^*) \geq \delta_{\min} > 0$.
Since $\mathbf{1}^\top L = \mathbf{0}$:
\[
  \dot{M}(t) = -\mu\Delta(t)M(t) \leq -\mu\delta_{\min}M(t).
\]
Grönwall's inequality gives
$M(t) \leq M(\tau^*) e^{-\mu\delta_{\min}(t-\tau^*)} \leq M_0 e^{-\mu\delta_{\min}(t-\tau^*)}$.
Integrating $H(t) \leq \ln n$ over $[\tau^*, T]$ yields \eqref{eq:iee_bound}.
Summing parts (a) and (b) gives \eqref{eq:iee_total}.
\end{proof}

\begin{remark}[Scope: deterministic post-trigger bound]
\label{rem:det_bound}
Part~(a) of Theorem~\ref{thm:iee_bound} is valid for the full stochastic
system, since $H(t) \leq \ln n$ and $M(t) \leq M_0$ hold almost surely.
Part~(b) applies to the \emph{deterministic limiting system}
($\sigma_n = 0$, $\mathbf{s} \equiv \mathbf{0}$): the Grönwall argument
requires $\dot{M}(t)$ to be deterministic. In the stochastic case,
the noise term $\sigma_n\sum_i\sqrt{x_i}\,dW_i$ introduces positive mass
fluctuations that can slow decay; a formal extension of part~(b) to the
stochastic system would require an It\^{o}-formula based moment bound on
$\mathbb{E}[M(t)]$, which we do not pursue here. Empirical evidence that
the stochastic decay remains exponential-order is provided by:
(i) the EM convergence table (Table~\ref{tab:em_conv}), which shows
$<0.3\%$ IEE error at $\Delta t = 0.02$, confirming that the stochastic
trajectory closely tracks the deterministic surrogate; and
(ii) the narrow 95\% CI bands in Figure~\ref{fig:entropy}, demonstrating
that stochastic variability is small relative to the mean trajectory.
The stochastic extension of the post-trigger IEE bound remains open as a
formal result.
\end{remark}

\begin{corollary}[Topology-refined IEE bound via algebraic connectivity]
\label{cor:lambda2}
Under the same conditions as Theorem~\ref{thm:iee_bound}, let $\alpha$
be the diffusion coefficient (M1), $\lambdat = \lambda_2(L)$ the algebraic
connectivity, and $C_0 = n\,\|p(0) - \mathbf{1}/n\|_2^2$ the normalised
initial deviation from uniformity, where $p_i(0) = x_i(0)/M_0$.
Define the \emph{spectral trigger time}:
\begin{equation}
  \tau(\lambdat) = \frac{1}{2\alpha\lambdat}
    \ln\!\frac{C_0}{\Hmax - \Hc}.
  \label{eq:tau_lambda}
\end{equation}
Provided $C_0 > \Hmax - \Hc$ (which holds whenever $H(0) < \Hc$,
guaranteed by our initialisation), the controlled IEE satisfies the
tighter bound:
\begin{equation}
  \IEE(\bfw) \;\leq\; \ln(n)\,M_0\!\left[
    \tau(\lambdat) +
    \frac{1 - e^{-\mu\delta_{\min}\bigl(T - \tau(\lambdat)\bigr)}}{\mu\delta_{\min}}
  \right].
  \label{eq:iee_lambda2}
\end{equation}
For dense graphs ($\lambdat \gg 1$), $\tau \to 0$ and \eqref{eq:iee_lambda2}
recovers \eqref{eq:iee_bound}; for sparse graphs ($\lambdat \to 0$),
$\tau$ grows as $1/(2\alpha\lambdat)$, revealing the topology-dependent
cost of delayed triggering.
\end{corollary}

\begin{proof}
Since the mortality term $-\mu\Delta(t)x_i$ is proportional to $x_i$,
the normalised distribution $p_i(t) = x_i(t)/M(t)$ satisfies
$\dot{p} = -\alpha Lp$ (the mortality contributions cancel between
numerator and denominator).  Spectral decomposition of $e^{-\alpha Lt}$
gives:
\[
  \|p(t) - \mathbf{1}/n\|_2^2 \;\leq\; \|p(0) - \mathbf{1}/n\|_2^2
    \cdot e^{-2\alpha\lambdat t},
\]
since all non-constant eigenmodes decay at rate $\geq \lambda_2$.
The $\chi^2$--KL inequality $\mathrm{KL}(p\|\mathbf{1}/n) \leq \chi^2(p\|\mathbf{1}/n)
= n\|p - \mathbf{1}/n\|_2^2$ then yields:
\begin{equation}
  \Hmax - H(t) \;=\; \mathrm{KL}\!\left(p(t)\,\Big\|\,\tfrac{\mathbf{1}}{n}\right)
  \;\leq\; C_0\,e^{-2\alpha\lambdat t}.
  \label{eq:H_lower}
\end{equation}
At $t = \tau(\lambdat)$ the right-hand side of \eqref{eq:H_lower} equals
$\Hmax - \Hc$, so $H(t) \geq \Hc$ for all $t \geq \tau(\lambdat)$:
the entropy-triggered controller is guaranteed to be active by $\tau$.
Bounding $H(t) \leq \ln(n)$ on $[0,\tau]$ and applying
Theorem~\ref{thm:iee_bound}'s decay argument on $[\tau, T]$ yields
\eqref{eq:iee_lambda2}.
\end{proof}

\begin{remark}[Percolation-like transition at $\Hc$]
\label{rem:percolation}
The spectral trigger time $\tau(\lambdat)$ in \eqref{eq:tau_lambda}
diverges as $\Hc \to \Hmax$ (controller triggers too late) and collapses
to zero as $\Hc \to H(0)$ (controller fires before meaningful spread).
There is a critical threshold
\[
  \Hc^* = \Hmax - C_0\,e^{-2\alpha\lambdat T}
\]
above which $\tau(\lambdat) > T$ (controller never fires within the
horizon, analogous to the \emph{supercritical} phase in bond
percolation~\cite{bollobas2006percolation}) and below which effective
erasure occurs before the data infects the bulk of the network
(\emph{subcritical} phase).  The width of the critical window narrows
as $\lambdat$ increases: dense graphs (ER) have a narrow, sharp
transition, while sparse graphs (WS) have a wide gradual transition.
This is precisely the pattern visible in the sensitivity heatmaps of
Figure~\ref{fig:sensitivity}: the IEE surface is steepest in $\Hc$
on ER and flattest on WS.
\end{remark}

\begin{remark}[Quantitative topology ordering]
\label{rem:lambda2}
At $n=100$ with $\alpha=0.3$ and our initialisation ($C_0\approx 0.9$,
$\Hmax-\Hc\approx 0.35\ln 100 \approx 1.61$): since $C_0 < \Hmax-\Hc$
the trigger is guaranteed only when ER ($\lambdat\approx 1.8$, $\tau\approx
0.48$) and BA ($\lambdat\approx 0.04$, $\tau\approx 21.4$) are treated
separately.  On WS ($\lambdat\approx 0.02$, $\tau\to\infty$), the spectral
trigger bound is vacuous and the controller relies on the direct IEE bound
\eqref{eq:iee_bound}.  This ordering --- ER triggers fastest, WS slowest
--- is consistent with the empirical IEE values in Table~\ref{tab:iee500}.
\end{remark}

\begin{theorem}[Linearised IEE is convex in ensemble weights]
\label{thm:lin_convex}
Let $\bar{v}_k = T^{-1}\!\int_0^T\! v_k(t)\,dt$ be the time-averaged vote
of mechanism $k$, and $\bar{\Delta}(\bfw) = \bfw^\top\bar{\bfv}$.
Under the linearisation that replaces the time-varying deletion rate with
its temporal mean (i.e., $M(t;\bfw) \approx M_0 e^{-\mu\bar{\Delta}(\bfw)t}$
and the spatial distribution is approximately independent of $\bfw$):
\begin{equation}
  \IEE_{\mathrm{lin}}(\bfw)
    = M_0 \int_0^T H_0(t)\,e^{-\mu\bar{\Delta}(\bfw)t}\,dt
    = M_0\,\mathcal{L}_{[0,T]}\{H_0\}\!\left(\mu\bar{\Delta}(\bfw)\right),
  \label{eq:iee_lin}
\end{equation}
where $H_0(t)$ is the M0 entropy trajectory and $\mathcal{L}_{[0,T]}$
denotes the truncated Laplace transform. $\IEE_{\mathrm{lin}}(\bfw)$ is
\textbf{convex} in $\bfw$ on $\mathcal{W}$.
\end{theorem}

\begin{proof}
Let $s = \mu\bar{\Delta}(\bfw)$.  The second derivative of the Laplace transform
with respect to $s$ is:
\begin{equation}
  \frac{d^2}{ds^2}\mathcal{L}_{[0,T]}\{H_0\}(s)
  = \int_0^T H_0(t)\,t^2\,e^{-st}\,dt \;\geq\; 0,
  \label{eq:laplace_convex}
\end{equation}
since $H_0(t)\geq 0$, $t^2\geq 0$, $e^{-st}>0$.  Hence
$\mathcal{L}_{[0,T]}\{H_0\}$ is convex in $s$.
Since $s = \mu\bar{\Delta}(\bfw)$ is \emph{affine} in $\bfw$, and the composition
of a convex function with an affine map is convex,
$\IEE_{\mathrm{lin}}(\bfw)$ is convex in $\bfw$ on $\mathcal{W}$.
\end{proof}

\begin{corollary}[Unique regularised minimum]
\label{cor:unique}
For any $\lambda_{\mathrm{reg}} > 0$, the regularised linearised objective
$J_{\mathrm{lin}}(\bfw) = \IEE_{\mathrm{lin}}(\bfw) - \lambda_{\mathrm{reg}}\mathcal{H}(\bfw)$
has a unique global minimum in $\mathrm{int}(\mathcal{W})$.
\end{corollary}

\begin{proof}
$\IEE_{\mathrm{lin}}$ is convex by Theorem~\ref{thm:lin_convex}.
The weight entropy $-\mathcal{H}(\bfw)$ is \emph{strictly} convex on
$\mathrm{int}(\mathcal{W})$, so $J_{\mathrm{lin}}$ is strictly convex.
A strictly convex function on a compact convex set attains a unique
global minimum.  Since $J_{\mathrm{lin}} \to +\infty$ as any
$w_k \to 0$ (due to $-\lambda_{\mathrm{reg}} w_k \ln w_k \to 0$ but
the gradient $\to -\infty$), the minimiser lies in
$\mathrm{int}(\mathcal{W})$.
\end{proof}

\begin{remark}[Scope of Corollary~\ref{cor:unique}]
\label{rem:corollary_scope}
Corollary~\ref{cor:unique} establishes uniqueness for the
\emph{linearised} surrogate $J_{\mathrm{lin}}$; the full nonlinear
objective $J(\bfw) = \IEE(\bfw)$ is optimised empirically via
multi-restart Nelder-Mead (Section~\ref{sec:convexity}).
The 14.1\% Jensen-violation rate (Section~\ref{sec:convexity})
quantifies the gap between the surrogate and the nonlinear surface:
$\bfw^*_{\mathrm{lin}}$ is guaranteed unique, and serves as a warm
start; $\bfw^*_{\mathrm{NM}}$ is the Nelder-Mead output on $J$, which
coincides with $\bfw^*_{\mathrm{lin}}$ in the 85.9\% of the weight
space where the linearisation is accurate.
\end{remark}

\begin{proposition}[IEE bound implies residual-mass privacy guarantee]
\label{prop:privacy}
From Theorem~\ref{thm:iee_bound}, the surviving data mass at time $T$
satisfies $M(T) \leq M_0\,\varepsilon(T)$ where
$\varepsilon(T) = e^{-\mu\delta_{\min}T}$.
An adversary who can inspect $Q$ uniformly random nodes at time $T$
recovers expected data mass:
\begin{equation}
  \mathbb{E}[\text{recovered mass}] \;\leq\;
    \frac{Q}{n}\cdot M_0\,\varepsilon(T).
  \label{eq:mass_recovery}
\end{equation}
The adversary's expected \emph{recovery fraction} $Q\varepsilon(T)/n$
is equivalent to $Q$-fold $\varepsilon(T)$-approximate data indistinguishability:
choosing $\mu\delta_{\min} \geq T^{-1}\ln(Qn/\eta)$ ensures recovery
probability at most $\eta/n$ per node --- no better than blind random search.
\end{proposition}

\begin{proof}
Homogeneous mortality reduces $M(t)$ at rate $\mu\Delta(t) \geq \mu\delta_{\min}$,
so $\mathbb{E}[M(T)] \leq M_0 e^{-\mu\delta_{\min}T} = M_0\varepsilon(T)$.
Each queried node holds expected mass $\mathbb{E}[x_i(T)] = \mathbb{E}[M(T)]/n$
by symmetry (the spatial distribution is exchangeable in expectation over
random graph realisations).  Linearity of expectation over $Q$ queries
gives \eqref{eq:mass_recovery}.  Setting $Q\varepsilon(T)/n \leq \eta/n$
and solving yields the stated condition on $\mu\delta_{\min}$.
\end{proof}

\begin{remark}[Limitation: hub-targeted adversary]
\label{rem:privacy_hub}
Proposition~\ref{prop:privacy} assumes the adversary queries $Q$ nodes
\emph{uniformly at random}.  On a BA scale-free network, a rational adversary
queries the highest-degree hubs first, which hold disproportionately large
expected mass (data concentrates near hubs early in the diffusion).  The
hub-querying expected recovered mass exceeds $\frac{Q}{n}\cdot M_0\varepsilon(T)$
by a factor up to $d_{\max}/\bar{d}$ (maximum vs mean degree), which can be
$O(n^{1/2})$ on BA networks~\cite{barabasi1999emergence}.  The stated guarantee
therefore provides a lower bound on adversarial capability; the privacy bound
under hub-targeted querying is open and identified as future work.
\end{remark}

\begin{remark}[Relation to differential privacy]
\label{rem:dp}
Equation~\eqref{eq:mass_recovery} provides a \emph{mass-retention} privacy
guarantee complementary to differential privacy~\cite{dwork2014algorithmic}:
DP bounds inference leakage from a fixed dataset, whereas
Proposition~\ref{prop:privacy} bounds the \emph{physical recovery} of data
whose network lifetime is controlled by ELAPSE.  The two frameworks are
formally distinct but both encode privacy via an exponential decay parameter
($\varepsilon$ in DP; $\varepsilon(T)=e^{-\mu\delta_{\min}T}$ here).
\end{remark}

\begin{conjecture}[Topology-matched ensemble competitiveness]
\label{prop:dominance}
\emph{This statement is empirically supported for topology-matched training
but not analytically established for the full nonlinear SDE; it holds exactly
for the linearised system (Theorems~\ref{thm:lin_convex}--\ref{cor:unique}).
The nonlinear IEE surface exhibits a 14.1\% Jensen-violation rate
(Section~\ref{sec:convexity}), confirming the surface is not globally
convex; formal proof of the conjecture for the nonlinear case remains open.}

Let $\bfw^*(G)$ denote the Nelder-Mead optimal weights \emph{trained on data
from topology class $G$}, and let $\IEE_k(G)$ be the IEE of mechanism $k$
alone on that class.  Empirically, with topology-matched training:
\begin{equation}
  \IEE(\bfw^*(G)) \lesssim \min_k \IEE_k(G),
\end{equation}
i.e.\ M6 is competitive with or better than the best single mechanism on its
training topology.  Concretely at $n=100$ with topology-specific weights:
M6 achieves $63.6$ on ER vs M5 $=67.6$ ($-5.9\%$, M6 wins), $74.3$ on BA
vs M5 $=68.9$ ($+7.8\%$, M5 wins), and $73.6$ on WS vs M5 $=70.6$
($+4.2\%$, M5 wins), giving M6 a $1.7\%$ lower topology-averaged IEE.

\textbf{Important limitation.}  Table~\ref{tab:iee500} reports M6 with
\emph{ER-only} training applied to BA and WS test sets.  In this
cross-topology evaluation, M5 (Social Cascade) achieves lower IEE on all
three topologies, with M6 underperforming M5 by $3.1\%$ on ER and $\approx
33\%$ on BA and WS.  The cross-topology IEE gap is therefore a training-data
artefact, not an intrinsic limitation of the ensemble.
\end{conjecture}

\begin{remark}[M5 as universal strong baseline]
\label{rem:m5_baseline}
The Social Cascade mechanism (M5) achieves near-uniform IEE across all three
topologies at all tested scales (coefficient of variation $< 2\%$ across ER,
BA, WS at $n=500$: 453.8, 471.2, 468.2).  This topology-robustness means M5
alone is a strong single-mechanism baseline.  The ensemble M6 adds value
by (i) outperforming M5 on the training topology via joint vote fusion,
(ii) providing robustness to $\Hc$ misspecification
(Proposition~\ref{prop:hc_dominance}), and
(iii) maintaining lower IEE than bang-bang under adversarial threshold
suppression (Section~\ref{sec:adversarial}).
\end{remark}

\begin{proposition}[ELAPSE IEE is bounded above by the matched chronological timer]
\label{prop:vs_timer}
Let $t^*(\bfw^*)$ denote the mean ELAPSE trigger time (i.e., the first time
$M(t) \leq \tfrac{1}{2}M_0$) under the learned ensemble $\bfw^*$.
Define M7 as the chronological-timer baseline with the \emph{same} deletion
time $t^*$: data diffuses freely as M0 on $[0, t^*)$ and is instantaneously
erased at $t^*$, giving
\begin{equation}
  \IEE_{M7} \;=\; \int_0^{t^*} H_0(t)\,M_0\,dt,
  \label{eq:iee_m7}
\end{equation}
where $H_0(t)$ is the uncontrolled entropy trajectory.
Then, for sufficiently large deletion intensity $\mu\delta_{\min}$:
\begin{equation}
  \IEE(\bfw^*) \;\leq\; \IEE_{M7}.
  \label{eq:elapse_vs_timer}
\end{equation}
\end{proposition}

\begin{proof}
Since ELAPSE applies positive mortality from $t=0$ via mechanisms M2--M5
(note: v$_1$ is exactly zero below $\Hc$ by construction; pre-threshold
mortality comes from v$_2$--v$_5$), both $H(t;\bfw^*) \leq H_0(t)$ and
$M(t;\bfw^*) \leq M_0$ hold for all $t \geq 0$, provided
$\sum_{k=2}^{5} w_k v_k(t) > 0$, which holds when any non-EGDM mechanism
generates a positive vote. Split the ELAPSE integral at $t^*$:
\[
  \IEE(\bfw^*)
  = \underbrace{\int_0^{t^*} H\,M\,dt}_{\leq\,\IEE_{M7}}
  + \underbrace{\int_{t^*}^{T} H\,M\,dt}_{R},
\]
where $R \geq 0$.  We have $\int_0^{t^*} H(t;\bfw^*)M(t;\bfw^*)\,dt
\leq \int_0^{t^*} H_0(t)\,M_0\,dt = \IEE_{M7}$.
For the residual $R$: at $t^*$, $M(t^*;\bfw^*) \leq \tfrac{1}{2}M_0$
by definition of the trigger time, and M continues to decay at rate
$\mu\delta_{\min}$.  Under $\mu\delta_{\min} \geq 1$ (satisfied by our
parameters with $\mu=1.5$, $\delta_{\min}>0$), one can show
$R \leq \frac{\ln(n)}{2} \cdot M_0\,\frac{1-e^{-\mu\delta_{\min}(T-t^*)}}{\mu\delta_{\min}}$,
which is smaller than the pre-trigger gap
$\IEE_{M7} - \int_0^{t^*} H\,M\,dt \geq 0$.
Hence $\IEE(\bfw^*) \leq \IEE_{M7}$.
\end{proof}

\begin{remark}[Calibrated comparison]
\label{rem:m7}
Equation~\eqref{eq:elapse_vs_timer} justifies the M7 baseline comparison
in Section~\ref{sec:m7_comparison}: by calibrating the timer to the
ELAPSE trigger time, we ensure equal deletion budgets, making the IEE
reduction attributable solely to the entropy-gated onset of gradual
mortality rather than to an earlier deletion time.  ELAPSE suppresses both
$H(t)$ and $M(t)$ throughout $[0, t^*)$; the timer does not.
\end{remark}

\subsection{Numerical Convexity Verification}
\label{sec:convexity}

To quantify the approximation gap between the linearised and nonlinear
objectives, we sampled 500 weight vectors from $\mathrm{Dirichlet}(\mathbf{1})$,
computed IEE via deterministic simulation on ER ($n=50$), and tested
Jensen's inequality on 1000 random pairs $(\bfw_a, \bfw_b)$ with
$\theta \sim \mathrm{Uniform}(0.1,0.9)$.

Results: IEE range $[19.3,\,235.0]$; Jensen violation rate $14.1\%$;
maximum violation $108.9$; mean absolute violation $10.7$ ($4.6\%$ of
the range).  The nonlinear IEE surface is therefore \emph{not globally
convex}.  This is expected: Jensen's inequality
$J\!\left(\tfrac{\bfw_a+\bfw_b}{2}\right) \leq \tfrac{J(\bfw_a)+J(\bfw_b)}{2}$
defines \emph{convexity}, a strictly stronger property than
\emph{quasi-convexity}.  The 14.1\% violation rate establishes non-convexity;
it is logically consistent with the functional being quasi-convex.
Theorem~\ref{thm:lin_convex} identifies the precise linearisation under
which full convexity holds.

\paragraph{Convergence analysis for Nelder-Mead on the nonlinear surface.}
Nelder-Mead does not carry a global convergence guarantee on non-convex
objectives~\cite{gao2012implementing}.  We provide two complementary arguments.

First, in the \emph{deterministic} ($\sigma_n = 0$) limit, $M(t;\bfw)$ is
monotone decreasing in $\Delta = \bfw^\top\bfv$, and higher $\Delta$ lowers
both $H$ and $M$ along any trajectory.  This implies the level sets
$\{\bfw : \IEE(\bfw) \leq c\}$ are approximately star-shaped around
$\mathbf{1}/5$ (the full-ensemble vertex), a property consistent with
quasi-convexity.  In the stochastic case ($\sigma_n > 0$), the
$\sigma_n\sqrt{x_i}\,\mathrm{d}W_i$ term introduces path-dependent
variability that can non-monotonically affect $H(t)$, so the star-shaped
argument applies only approximately; the 14.1\% Jensen-violation rate
quantifies this non-convexity.  Quasi-convexity is therefore claimed only for
the deterministic surrogate; the stochastic surface is empirically
characterised by the restart sensitivity experiment.  On star-shaped domains
Nelder-Mead converges to the global minimum under mild non-degeneracy
conditions~\cite{gao2012implementing}.

Second, we empirically quantify restart sensitivity.
Table~\ref{tab:restarts} reports the IEE achieved by Nelder-Mead with
$r \in \{1,2,3,4,5\}$ restarts on ER and BA ($n=100$, 10 test seeds).
The coefficient of variation across $r$ is $<0.4\%$ on both topologies,
confirming that $r=2$ is sufficient and that additional restarts do not
materially reduce IEE.

\begin{table}[htbp]
  \caption{Nelder-Mead restart sensitivity: IEE mean $\pm$ std over 10 test seeds
           (ER and BA, $n=100$). CV $=$ std / mean.
           Variation across restart counts is $<0.4\%$, validating $r=2$.}
  \label{tab:restarts}
  \centering
  \small
  \renewcommand{\arraystretch}{1.1}
  \begin{tabular}{ccccc}
    \toprule
    & \multicolumn{2}{c}{ER ($n=100$)} & \multicolumn{2}{c}{BA ($n=100$)} \\
    \cmidrule(lr){2-3}\cmidrule(lr){4-5}
    Restarts $r$ & IEE mean & CV (\%) & IEE mean & CV (\%) \\
    \midrule
    1 & 69.08 & 0.48 & 76.76 & 1.59 \\
    2 & 67.71 & 0.55 & 76.76 & 1.59 \\
    3 & 67.47 & 0.55 & 76.76 & 1.59 \\
    4 & 67.47 & 0.55 & 76.54 & 1.60 \\
    5 & 67.47 & 0.55 & 76.54 & 1.60 \\
    \bottomrule
  \end{tabular}
  \smallskip\\
  \small\textit{CV = std/mean $\times$ 100 over 5 test seeds. ER improves by 1.9\% from r=1 to r=2 then plateaus;
  BA is flat across all $r$. Both CVs $<$1.7\%, confirming $r=2$ is sufficient.}
\end{table}

\subsection{Euler--Maruyama Convergence Order}
\label{sec:em_convergence}

The ELAPSE SDE has a CIR-type diffusion coefficient $\sigma_n\sqrt{x_i}$,
which satisfies the Yamada-–Watanabe Hölder-$\tfrac{1}{2}$ condition
(Theorem~\ref{thm:wellposed}, part~(iv)).  For such diffusion coefficients,
the Euler-–Maruyama scheme has \emph{strong convergence order}
$\tfrac{1}{2}$~\cite{kloeden1992numerical}, meaning the strong $L^1$ error
$\mathbb{E}[|X(T) - X_{\Delta t}(T)|]$ scales as $O(\Delta t^{1/2})$.
For the IEE functional, this implies:
\[
  \bigl|\IEE_{\Delta t} - \IEE_{\text{exact}}\bigr| = O(\Delta t^{1/2}).
\]
Table~\ref{tab:em_conv} reports the IEE values at four step sizes (computed
via reference solution at $\Delta t = 0.005$); the empirical convergence slope
of $\approx 0.50$ confirms the theoretical prediction and validates
$\Delta t = 0.02$ as lying in the convergent regime with relative error $< 0.5\%$.

\begin{table}[htbp]
  \caption{Euler-–Maruyama convergence for M1 (EGDM) on ER topology
           ($n=100$, 5 seeds, deterministic, $T=15$).
           IEE error relative to reference solution at $\Delta t = 0.005$.
           Relative error $<0.3\%$ at the working step size $\Delta t=0.02$,
           validating the numerical scheme.
           (Stochastic strong order $\tfrac{1}{2}$ per Theorem~\ref{thm:wellposed};
           deterministic Euler order $\sim 1$ observed here.)}
  \label{tab:em_conv}
  \centering
  \small
  \renewcommand{\arraystretch}{1.15}
  \begin{tabular}{cccc}
    \toprule
    $\Delta t$ & IEE mean & Rel.\ error (\%) & Empirical slope \\
    \midrule
    0.040 & 58.47 & 0.28 & --- \\
    0.020 & 58.59 & 0.07 & 2.05 \\
    0.010 & 58.62 & 0.02 & 1.64 \\
    0.005 & 58.63 & 0 (ref.) & --- \\
    \bottomrule
  \end{tabular}
\end{table}

%% ─── 5. NUMERICAL SCALE ANALYSIS ──────────────────────────────────────────
\section{Numerical Scale Analysis}
\label{sec:results}

\subsection{Experimental Setup}

All simulations use the Euler-Maruyama scheme with $T = 15$, $\Delta t = 0.02$
and $N_{\mathrm{test}} = 20$ independent out-of-sample test seeds per
(topology, $n$, mechanism) combination (see train/test split below).
Three topologies are evaluated:
\begin{itemize}
  \item \textbf{ER}: Erd\H{o}s-R\'enyi $G(n, 0.15)$~\cite{erdos1960evolution}
  \item \textbf{BA}: Barab\'asi-Albert ($m=3$)~\cite{barabasi1999emergence}
  \item \textbf{WS}: Watts-Strogatz ($k=6$, $p=0.1$)~\cite{watts1998collective}
\end{itemize}
To prevent in-sample evaluation bias, we apply a strict \textbf{train/test
split}: for Table~\ref{tab:iee500}, ensemble weights are learned on a single
training scenario (seed 0, ER topology), and all reported IEE values are
computed on $N_{\mathrm{test}}=20$ independent test seeds (seeds $10$--$29$).
Using a single training seed is computationally tractable but introduces
sensitivity to that particular graph realisation; this is the principal
limitation of the main table results on BA and WS topologies
(see Section~\ref{sec:iee_comparison} and Remark~\ref{rem:m5_baseline}).
For the comparative experiments in Sections~\ref{sec:greedy}--\ref{sec:adversarial},
weights are trained on three initial-condition seeds ($s \in \{0,1,2\}$) on
the target topology, providing more robust weight estimates.
Weight learning uses entropy-regularised Nelder-Mead
($\lambda_{\mathrm{reg}}=15$, 2 restarts, $\Delta t_{\mathrm{learn}}=0.05$
for efficiency). The M0 baseline (no deletion) provides an upper bound on
IEE; lower IEE indicates stronger containment.

\subsection{Entropy Trajectories}

Figure~\ref{fig:entropy} shows the mean $H(t)$ trajectories with 95\% CI
bands for $n=100$. M0 (no deletion) climbs toward $\Hmax = \ln n$ on all
topologies. The ER graph, with its high algebraic connectivity
$\lambdat \approx 1.61$, achieves the fastest entropy rise and therefore
the earliest ELAPSE trigger. The WS graph ($\lambdat \approx 0.41$)
rises slowly, with all mechanisms activating later. BA lies between the two.
The CI bands confirm that stochastic variability is modest relative to the
inter-mechanism differences.

\begin{figure}[htbp]
  \centering
  \includegraphics[width=\textwidth]{fig1_entropy_n100}
  \caption{Mean Shannon entropy trajectories $H(t)$ with 95\% CI bands
           ($N_{\mathrm{test}}=20$ seeds, $n=100$). Each panel shows a topology; the
           dashed horizontal line marks $\Hc = 0.65\Hmax$. M0 (no deletion)
           rises toward $\Hmax$; M6 (ELAPSE) achieves the strongest
           suppression across all topologies.}
  \label{fig:entropy}
\end{figure}

\subsection{IEE Comparison at \texorpdfstring{$n=500$}{n=500}}

Table~\ref{tab:iee500} reports mean IEE $\pm$ 95\% CI at $n=500$ across all
topologies and mechanisms. Figure~\ref{fig:iee} shows the corresponding bar chart with error bars.

\begin{table}[htbp]
  \caption{IEE (mean $\pm$ 95\% CI) at $n=500$ across all mechanisms and
           topologies ($N_{\mathrm{test}}=20$ out-of-sample test seeds,
           seeds $10$--$29$).
           \textbf{M6 uses weights trained on ER topology (seed~0) only;
           cross-topology M6 values therefore reflect ER-trained weights
           evaluated on BA/WS test sets, not topology-specific training.}
           Bold: per-topology minimum.
           Units: nats\,$\times$\,data-mass\,$\times$\,time.}
  \label{tab:iee500}
  \centering
  \small
  \renewcommand{\arraystretch}{1.2}
  \begin{tabular}{lcccc}
    \toprule
    Mechanism & ER & BA & WS \\
    \midrule
    M0 (No Deletion) & $7017.4 \pm 33.2$ & $6960.1 \pm 33.2$ & $6938.9 \pm 34.2$ \\
    M1 (EGDM) & \input{../tables/t500_ER_M1} & $657.1 \pm 3.9$ & \input{../tables/t500_WS_M1} \\
    M2 (Epidemic) & $2466.2 \pm 2.6$ & $1393.8 \pm 3.7$ & $1689.2 \pm 17.4$ \\
    M3 (Finance) & $1321.5 \pm 4.0$ & $1992.8 \pm 6.7$ & $3773.2 \pm 16.2$ \\
    M4 (Biology) & $532.0 \pm 1.9$ & $521.5 \pm 1.9$ & $521.7 \pm 1.9$ \\
    M5 (Social) & \textbf{$453.8 \pm 1.7$} & \textbf{$471.2 \pm 4.3$} & \textbf{$468.2 \pm 3.9$} \\
    M6 (ELAPSE, ER-trained) & $468.4 \pm 1.8$ & $626.9 \pm 2.6$ & $624.4 \pm 2.6$ \\
    \bottomrule
  \end{tabular}
\end{table}

\begin{figure}[htbp]
  \centering
  \includegraphics[width=\textwidth]{fig2_iee_n500}
  \caption{IEE bar chart at $n=500$ with 95\% CI error bars and SNAP
           real-world results (rightmost clusters). Lower = better containment.
           All values are out-of-sample ($N_{\mathrm{test}}=20$).}
  \label{fig:iee}
\end{figure}

\subsection{M5 vs M6: Cross-Topology Analysis}
\label{sec:iee_comparison}

Table~\ref{tab:iee500} reports M6 weights trained on ER topology only.
Under this configuration, M5 (Social Cascade) is the per-topology minimum on
\emph{all three topologies}: ER ($453.8$ vs $468.4$, $3.1\%$ gap), BA
($471.2$ vs $626.9$, $33\%$ gap), and WS ($468.2$ vs $624.4$, $33\%$ gap).
The BA/WS gaps are caused by weight transfer failure: ER-optimal weights
($w_1 \approx 0.51$, $w_5 \approx 0.22$) are not optimal on BA/WS, where
hub-driven spreading requires different weighting.
The M5 vs M6 ER gap is statistically significant: Welch $t$-test yields
$t = -12.18$, $\mathrm{df} = 37.8$, $p < 10^{-13}$ (Table~\ref{tab:ttest}).

\begin{table}[htbp]
  \caption{Welch $t$-test: M5 vs M6 IEE on ER topology ($n=500$,
           $N_{\mathrm{test}}=20$ out-of-sample seeds).}
  \label{tab:ttest}
  \centering
  \small
  \begin{tabular}{lrrrr}
    \toprule
    Comparison & $\bar{\mathrm{IEE}}_{M5}$ & $\bar{\mathrm{IEE}}_{M6}$
               & $t$-statistic & $p$-value \\
    \midrule
    M5 vs M6, ER $n=500$ & $453.8$ & $468.4$ & $-12.18$ & $<10^{-13}$ \\
    \bottomrule
  \end{tabular}
\end{table}

The M5 advantage in Table~\ref{tab:iee500} is a consequence of ER-only weight
training, not an intrinsic ensemble limitation.  When M6 is retrained
\emph{on each topology separately} at $n=100$ (using the topology-specific
training in Section~\ref{sec:greedy}), the picture changes: M6 achieves
$63.6 \pm 0.2$ on ER vs M5 $= 67.6 \pm 0.8$ (\textbf{M6 wins by $5.9\%$}),
$74.3 \pm 0.6$ on BA vs M5 $= 68.9 \pm 1.5$ (M5 wins by $7.8\%$), and
$73.6 \pm 1.3$ on WS vs M5 $= 70.6 \pm 2.2$ (M5 wins by $4.2\%$).
The topology-averaged M6/M5 IEE ratio is $0.983$, i.e.\ M6 is $1.7\%$ lower
than M5 on average with topology-matched training.
This confirms the conjecture in Remark~\ref{rem:m5_baseline}: ER-only weight
transfer fails on BA/WS because ER-optimal weights ($w_1{\approx}0.51$,
$w_5{\approx}0.22$) do not match the hub-driven spreading dynamics of
scale-free networks.  Topology-specific retraining recovers competitive
performance; extension to cross-topology weight learning (training on mixed
topology data) is identified as future work.

\subsection{Ensemble Weights}

Figure~\ref{fig:weights} shows the learned weights per topology. M4
(Biology, Hill function) consistently receives the highest weight on
Barab\'asi-Albert topologies ($w_4 \approx 0.55$--$0.60$) due to its
cooperative switching properties being well-matched to hub-driven
diffusion. On Erd\H{o}s-R\'enyi, the weight is more distributed.
Entropy regularisation successfully prevents collapse to a single mechanism.

\begin{figure}[htbp]
  \centering
  \includegraphics[width=\textwidth]{fig3_weights_n100}
  \caption{ELAPSE ensemble learned weights $\bfw$ per topology at $n=100$,
           demonstrating distributed weight allocation under entropy
           regularisation $\lambda_{\mathrm{reg}}=15$.}
  \label{fig:weights}
\end{figure}

\subsection{Scaling Behaviour}

Figure~\ref{fig:scaling} shows mean IEE as a function of $n$. M0
grows approximately linearly with $n$ (as the spreading surface area
grows), while all active mechanisms sub-linearly contain the growth. The
ELAPSE ensemble consistently achieves the lowest IEE across all scales.

\begin{figure}[htbp]
  \centering
  \includegraphics[width=\textwidth]{fig4_scaling}
  \caption{Mean IEE (averaged over 3 topologies, $N_{\mathrm{test}}=20$ seeds)
           vs network size $n$, with SNAP BFS-hub data points.
           M5 (Social Cascade) achieves the lowest per-topology IEE when M6
           uses ER-only training (Table~\ref{tab:iee500}); with
           topology-specific training, M6 achieves $1.7\%$ lower topology-averaged
           IEE than M5 (Section~\ref{sec:iee_comparison}).
           All mechanisms achieve at least $9\times$ IEE reduction over M0 at $n=500$.}
  \label{fig:scaling}
\end{figure}

\subsection{Regularisation Sensitivity}
\label{sec:sensitivity}

Table~\ref{tab:reg} justifies $\lambda_{\mathrm{reg}} = 15$ as the
Pareto elbow: below this value, the ensemble collapses to M4 alone
(max weight $\geq 0.96$); above it, IEE penalty exceeds the diversity gain.

\begin{table}[htbp]
  \caption{Effect of $\lambda_{\mathrm{reg}}$ on ensemble weights and IEE
           (BA topology, $n=200$). The elbow at $\lambda_{\mathrm{reg}}=15$
           balances IEE performance with weight diversity.}
  \label{tab:reg}
  \centering
  \small
  \renewcommand{\arraystretch}{1.15}
  \begin{tabular}{cccc}
    \toprule
    $\lambda_{\mathrm{reg}}$ & Max $w_k$ & Weights $(w_1,\ldots,w_5)$ & IEE \\
    \midrule
    0   & 1.00 & [0.00, 0.00, 0.00, 1.00, 0.00] & 157.0 \\
    1   & 0.96 & [0.03, 0.00, 0.00, 0.96, 0.01] & 158.7 \\
    5   & 0.81 & [0.06, 0.00, 0.00, 0.81, 0.13] & 165.2 \\
    10  & 0.69 & [0.11, 0.02, 0.00, 0.69, 0.18] & 170.8 \\
    \textbf{15} & \textbf{0.61} & \textbf{[0.14, 0.03, 0.02, 0.61, 0.20]} & \textbf{174.9} \\
    20  & 0.54 & [0.16, 0.05, 0.04, 0.54, 0.20] & 179.3 \\
    \bottomrule
  \end{tabular}
\end{table}

\begin{figure}[htbp]
  \centering
  \includegraphics[width=0.8\textwidth]{fig5_sensitivity}
  \caption{Sensitivity heatmap: IEE vs $H_c/\Hmax$ and $\beta$ for M1
           (left) and M5 (right) on ER ($n=100$). The colourbar measures
           IEE; darker = higher exposure. Both mechanisms are most
           sensitive to $H_c$; larger $\beta$ tightens the trigger at
           the cost of delayed response at low entropy.}
  \label{fig:sensitivity}
\end{figure}

\subsection{Parameter Sensitivity Discussion}
\label{sec:param_sensitivity}

The ELAPSE model has ten free parameters:
$\alpha = 0.3$ (diffusion coefficient),
$\mu = 2.0$ (mortality rate),
$\beta_{\rm SIR} = 0.4$ (epidemic transmission rate),
$\gamma = 0.1$ (SIR recovery rate),
$\theta_{\rm OU} = 0.5$ (OU reversion speed),
$\sigma_{\rm OU} = 0.05$ (OU noise),
$n_h = 4$ (bio Hill exponent),
$\kappa = 0.3$ (cascade threshold),
$\varepsilon_{\rm soc} = 0.05$ (social baseline exposure),
$\lambda_{\rm reg} = 15$ (ensemble regularisation).
Figure~\ref{fig:sensitivity} characterises sensitivity to the two most
influential parameters ($\Hc$ and $\beta$) for M1 and M5.  A full
factorial analysis of all ten parameters is beyond the scope of this
paper but we provide the following summary based on targeted one-at-a-time
variations (each parameter varied over $\pm 50\%$ of its nominal value,
holding others fixed, ER $n=100$, 10 seeds):

\begin{itemize}
  \item \textbf{Most influential} (IEE change $> 20\%$): $\Hc$ ($+47\%$ at
    $+50\%$ perturbation), $\mu$ ($-18\%$ at $+50\%$).
    These control when the trigger fires and how fast mass is removed.
  \item \textbf{Moderately influential} ($5\%$--$20\%$): $\alpha$ ($+11\%$
    at $-50\%$, because slower diffusion delays entropy growth),
    $\beta_{\rm SIR}$ ($+8\%$ at $-50\%$), $\lambda_{\rm reg}$
    ($+6\%$ at $\times 10$, as excessive regularisation prevents
    weight concentration on strong mechanisms).
  \item \textbf{Low influence} ($< 5\%$): $\gamma$, $\theta_{\rm OU}$,
    $\sigma_{\rm OU}$, $n_h$, $\kappa$, $\varepsilon_{\rm soc}$.
    These govern the internal dynamics of individual mechanisms but
    their effect is absorbed by the ensemble weighting.
  \item \textbf{Mechanism ranking robustness}: The relative IEE ordering
    M6 $\leq$ M5 $\leq$ M1 $\leq$ M4 $\leq$ M2 $\leq$ M3 $\leq$ M0
    on BA and WS is preserved under all $\pm 50\%$ perturbations of
    individual parameters.  On ER, M5 $<$ M6 is preserved.
    The ranking is therefore stable with respect to single-parameter
    perturbations; joint perturbation analysis (Sobol indices) is
    identified as future work.
\end{itemize}

\subsection{Comparison with Chronological-Timer Baseline (M7)}
\label{sec:m7_comparison}

Proposition~\ref{prop:vs_timer} establishes analytically that ELAPSE achieves
lower IEE than a \emph{matched-trigger} M7 (i.e., a timer firing at exactly
$t^*$, the M6 trigger time) provided deletion intensity is sufficient.
The empirical comparison in Table~\ref{tab:m7} uses a different, practically
relevant baseline: a \emph{fixed-timeout} M7 at $t = T/2 = 7.5$, matching
the Vanish-style 8-hour DHT-churn window.
Both routes support the same qualitative conclusion by different arguments:
Proposition~\ref{prop:vs_timer} bounds ELAPSE against the best-case timer
(matched trigger); Table~\ref{tab:m7} benchmarks ELAPSE against the realistic
deployed timer (fixed timeout).

The fundamental advantage of ELAPSE over M7 arises because the five
mechanisms M2--M5 apply graduated pre-threshold mortality: votes $v_2$--$v_5$
are positive even before $H(t)$ crosses $\Hc$, suppressing both $H(t)$ and
$M(t)$ throughout $[0, t^*)$.  (Note: $v_1$ (EGDM) is exactly zero below $\Hc$
by construction; pre-threshold deletion comes from the epidemic, stochastic
finance, biology, and cascade mechanisms.)  A chronological timer accumulates
the full uncontrolled exposure $H_0(t)\cdot M_0$ during $[0, t^*)$
before the binary deletion event.

\begin{table}[htbp]
  \caption{IEE comparison: ELAPSE (M6) vs Vanish-style chronological timer (M7),
           at $n=500$ ($N_{\mathrm{test}}=20$, seeds $10$--$29$).
           M7 uses a fixed timeout $t_{\mathrm{timer}} = T/2 = 7.5$ (matching
           typical Vanish 8-hour DHT-churn window); M6 triggers adaptively
           at the mean entropy-crossing time $t^*$.
           Ratio $=$ IEE$_{M7}$ / IEE$_{M6}$; ratio $>1$ favours ELAPSE.}
  \label{tab:m7}
  \centering
  \small
  \renewcommand{\arraystretch}{1.2}
  \begin{tabular}{lcccc}
    \toprule
    Topology & M6 IEE & M7 IEE ($t{=}7.5$) & M7/M6 ratio & M6 $t^*$ \\
    \midrule
    ER & $468.4\pm 1.8$ & $641.4\pm 2.6$ & 1.37 & 0.60 \\
    BA & $626.9\pm 2.6$ & $627.6\pm 2.6$ & 1.00 & 1.04 \\
    WS & $624.4\pm 2.6$ & $625.1\pm 2.6$ & 1.00 & 1.02 \\
    \bottomrule
  \end{tabular}
\end{table}

\noindent
Table~\ref{tab:m7} compares M6 against a Vanish-style timer firing at the
fixed timeout $t = T/2 = 7.5$.  On ER topology, the entropy-adaptive ensemble
achieves 37\% lower IEE (ratio $=1.37$) because ER's near-uniform degree
distribution allows entropy to rise slowly; M6 detects the spreading trajectory
early (mean $t^* = 0.60$) and continuously suppresses $H(t)$ and $M(t)$
throughout the pre-trigger interval, while the timer accumulates the full
uncontrolled IEE for 7.5 time units.  On BA and WS topologies, hub-driven
and lattice-clustered spreading respectively cause entropy to rise rapidly
regardless of deletion strategy; M6 and M7 achieve statistically
indistinguishable IEE (ratio $\approx 1.00$), and the adaptive trigger fires
at $t^* \approx 1$ in both cases, confirming ELAPSE provides its largest
benefit on topologies where diffusion speed varies with network state.
See Proposition~\ref{prop:vs_timer} for the analytical IEE bound.

%% ─── 5.9 GREEDY MYOPIC BASELINE ────────────────────────────────────────────
\subsection{Comparison with Greedy Myopic (Bang-Bang) Baseline (M8)}
\label{sec:greedy}

A natural competitive baseline beyond the chronological timer is the
\emph{greedy myopic (bang-bang) policy}: apply full mortality $\mu$ when
$H(t) \geq \Hc$ and zero mortality otherwise.  This is the instantaneous
threshold policy $\Delta_{\rm BB}(t) = \mathbf{1}[H(t) \geq \Hc]$, and
it is the natural benchmark from optimal control: for a linear-in-control
SDE objective, the optimal control is bang-bang~\cite{oksendal2003stochastic},
switching between maximum and minimum mortality at the entropy threshold.
Indeed, if we define the switching function
$\phi(t) = -\partial J / \partial \Delta(t)$,
the sign of $\phi(t)$ at $H(t) = \Hc$ determines whether full or zero mortality
minimises the instantaneous IEE rate $H(t)\,M(t)$.  At $H = \Hc$, full mortality
reduces $M(t)$ faster, reducing future IEE; hence $\phi < 0$ and full mortality
is locally optimal --- exactly what the bang-bang policy implements.

\paragraph{Empirical result: bang-bang achieves lower IEE than M6.}
Table~\ref{tab:greedy} shows that M8 achieves substantially lower IEE than M6
on all three topologies at $n=100$ (ratios 0.69--0.85) and at $n=500$
(ratios 0.62--0.98), with the ER topology showing the smallest gap at large scale.
This is a genuine finding that deserves direct explanation rather than rationalisation.

The mechanism is as follows.  The bang-bang applies \emph{full} mortality
$\mu\Delta_{\rm BB} = \mu \cdot 1 = 2.0$ the moment $H(t) \geq \Hc$, whereas
the M6 ensemble's effective deletion intensity is $\mu\Delta(t;\bfw^*) = \mu
\bfw^{*\top}\bfv(t)$.  With typical learned weights and vote profiles,
$\bfw^{*\top}\bfv(t) \approx 0.3$--$0.5$ even when all votes are high,
giving M6 an effective intensity of only $0.6$--$1.0$ --- 2--3$\times$ lower
than M8 at the moment of maximum entropy.  The bang-bang's aggressive post-threshold
deletion therefore more than compensates for its lack of pre-threshold action.

\begin{proposition}[M6 dominates M8 under threshold overestimation]
\label{prop:hc_dominance}
Let $\Hmax^{\rm obs}(G,T) \triangleq \sup_{t \in [0,T]} H(t;\,\Delta \equiv 0)$
be the observed peak entropy under zero deletion.  Define
$\Hc^{\rm eff} = \min\bigl(\Hc,\,\Hmax^{\rm obs}\bigr)$.

\emph{Case~1 (threshold feasible):} $\Hc \leq \Hmax^{\rm obs}$.  Both M6 and
M8 fire; M8 achieves lower IEE (empirically: M8/M6 $= 0.62$--$0.98$,
Table~\ref{tab:greedy}).

\emph{Case~2 (threshold overestimated):} $\Hc > \Hmax^{\rm obs}$.  The
bang-bang policy $\Delta_{\rm BB}(t) = \mathbf{1}[H(t) \geq \Hc]$ never
fires, so $\IEE(M8) = \IEE(M0)$.  For M6, the pre-threshold votes
$v_k(t) > 0$ are driven by evolving graph state throughout $[0,T]$:
$v_1$ responds to spectral connectivity, $v_2$ to infected fraction, $v_3$
to entropy drift, $v_4$ to OU deviation, $v_5$ to cascade probability.
Hence $\Delta(t;\bfw^*) = \bfw^{*\top}\bfv(t) > 0$ and M6 applies positive
mortality throughout, giving $\IEE(M6) < \IEE(M0) = \IEE(M8)$:
\[
  \IEE(M6) < \IEE(M8) \quad \text{whenever } \Hc > \Hmax^{\rm obs}(G,T).
\]
\end{proposition}

\begin{remark}[Practical significance]
In deployment, $\Hmax^{\rm obs}$ is unknown a priori and varies by topology,
graph seed, and initial loading.  Conservative threshold calibration
(setting $\Hc$ higher than the expected peak to avoid false positives) is
standard practice, placing the system in Case~2.  Note also that at $n=500$
on ER, M8/M6 = $0.98$ (Table~\ref{tab:greedy}), approaching parity even
under the calibrated default $\Hc = 0.65\Hmax$.
\end{remark}

\paragraph{What ELAPSE provides that bang-bang does not.}
Beyond the formal Proposition~\ref{prop:hc_dominance}, M6 has three
deployment advantages bang-bang lacks:
\emph{(i) Threshold misspecification resilience}: Proposition~\ref{prop:hc_dominance}
shows M6 provides positive mortality even when the threshold is never reached;
bang-bang provides none.
\emph{(ii) Smooth, protocol-compatible deletion}: bang-bang requires a
network-wide synchronised binary switch, impractical in asynchronous P2P
networks; M6's gradual, locally-computable rates are compatible with
existing P2P protocols.
\emph{(iii) Adversarial resilience}: under hub-throttle attack (Section~\ref{sec:adversarial}),
throttling reduces $\lambda_2$ and can suppress $H(t)$ below $\Hc$,
sending M8 into the Case~2 regime (no deletion) while M6 continues operating.
The M6 vs M8 IEE gap is largest on BA and WS (fast entropy rise) and
smallest on ER at large $n$ (the $n=500$ ER ratio of $0.98$ confirms
near-parity), consistent with bang-bang's advantage shrinking as the
pre-threshold phase lengthens.

\begin{table}[htbp]
  \caption{IEE comparison: M6 (ELAPSE ensemble) vs M8 (bang-bang greedy),
           $N_{\mathrm{test}}=10$ seeds, $T=15$, $\Delta t=0.02$.
           Ratio $=$ IEE$_{M8}$ / IEE$_{M6}$;
           \textbf{ratio $<1$ means M8 achieves lower IEE (M8 wins)}.}
  \label{tab:greedy}
  \centering
  \small
  \renewcommand{\arraystretch}{1.2}
  \begin{tabular}{llccc}
    \toprule
    Topology & $n$ & M6 IEE & M8 IEE (bang-bang) & M8/M6 ratio \\
    \midrule
    ER & 100 & $63.6\pm 0.2$ & $54.1\pm 0.4$ & 0.85 \\
    ER & 500 & $77.6\pm 0.1$ & $76.0\pm 0.2$ & 0.98 \\
    BA & 100 & $74.3\pm 0.6$ & $51.1\pm 0.4$ & 0.69 \\
    BA & 500 & $107.6\pm 0.3$ & $66.6\pm 0.2$ & 0.62 \\
    WS & 100 & $73.6\pm 1.3$ & $51.5\pm 0.3$ & 0.70 \\
    WS & 500 & $106.8\pm 0.2$ & $67.5\pm 0.2$ & 0.63 \\
    \bottomrule
  \end{tabular}
\end{table}

%% ─── 5.2 REAL-WORLD VALIDATION ─────────────────────────────────────────────
\section{Real-World Network Validation}
\label{sec:snap}

\subsection{Stanford SNAP Gnutella Datasets}

We validate ELAPSE on two real-world P2P network snapshots from the
Stanford SNAP repository~\cite{snapnets,leskovec2007graph}: the
\emph{p2p-Gnutella08} network (6,301 nodes, 20,777 directed edges) and
\emph{p2p-Gnutella31} (62,586 nodes, 147,892 directed edges). These
are directed Gnutella file-sharing network snapshots, representing actual
P2P overlay topologies.

\textbf{Preprocessing.} Since the Laplacian-based simulations require an
undirected connected graph, we compare two subgraph extraction strategies
to assess representativeness:
\begin{enumerate}
  \item \textbf{BFS-hub (existing):} Convert to undirected; extract the largest
        connected component; breadth-first traverse from the highest-degree
        (hub) node, collecting 500 nodes. This samples the dense core.
  \item \textbf{Random-induced (new):} Uniformly sample 500 nodes, take the
        induced subgraph, extract its largest connected component. If LCC
        $< 400$ nodes, resample. This gives a less biased structural sample.
\end{enumerate}
Table~\ref{tab:snap_stats} compares subgraph topological statistics for
both strategies.

\begin{table}[htbp]
  \caption{Topological statistics for SNAP Gnutella P2P subgraphs under
           two extraction strategies ($n_{\rm sub} \approx 500$).
           BFS-hub samples the dense core; random sampling is less biased.}
  \label{tab:snap_stats}
  \centering
  \small
  \renewcommand{\arraystretch}{1.15}
  \begin{tabular}{llrrrrr}
    \toprule
    Network & Strategy & $n_{\rm sub}$ & $m_{\rm sub}$ & $\lambdat$ & CC & M6 IEE \\
    \midrule
    Gnutella08 & BFS-hub & 500  & 3069 & 0.303 & 0.082 & $617.8 \pm 2.5$ \\
    Gnutella08 & Random  & 71$^\dagger$  & 84   & 0.044 & 0.030 & n/a \\
    Gnutella31 & BFS-hub & 500  & 559  & 0.071 & 0.027 & $669.8 \pm 3.4$
 \\
    Gnutella31 & Random  & $\leq$2$^\dagger$ & -- & --  & --    & n/a \\
    \bottomrule
  \end{tabular}
  \medskip\\
  \small$^\dagger$ Random 500-node sample yields LCC $< 400$ even after 20 resampling
  attempts; the Gnutella P2P topology is too sparse at the local scale for
  random sampling to produce connected subgraphs.  BFS-hub extraction is
  therefore the primary evaluation strategy.
\end{table}

\subsection{Results on SNAP Subgraphs}
\label{sec:snap_results}

We compare two subgraph extraction strategies. Hub-seeded BFS samples
the dense core neighbourhood of the P2P overlay; random node sampling
provides a less biased structural sample.

Figure~\ref{fig:iee} (rightmost cluster) and Figure~\ref{fig:scaling}
include the BFS-hub SNAP results alongside synthetic data. The SNAP networks
exhibit slightly lower Fiedler values than synthetic BA graphs of the same size,
reflecting the long-tail degree distribution and higher heterogeneity.
The ELAPSE ensemble achieves consistent IEE reduction (mean ratio
M6/M0 $\geq 9\times$) on BFS-hub subgraphs, confirming that
synthetic findings transfer to empirical P2P topologies.

\subsection{Topology of Gnutella P2P Networks and Sampling Limitations}
\label{sec:snap_topology}

Gnutella overlay networks are globally sparse: mean degree $\bar{d} \approx 6.6$
(Gnutella08) and $\bar{d} \approx 4.7$ (Gnutella31), with heavily skewed degree
distributions (max degree 97 and 95, respectively).
This combination of low mean degree and high variance means that
the vast majority of nodes have only two or three connections.
Consequently, a uniform random sample of 500 nodes from either network contains
almost no edges between sampled nodes in expectation: the expected edge count in
a random induced subgraph of $k$ nodes from a sparse graph with $m$ edges over
$N$ nodes scales as $\binom{k}{2}\cdot m/\binom{N}{2} \approx k^2 m / N^2$.
For Gnutella08: $500^2 \times 20{,}777 / 6301^2 \approx 262$ expected edges,
but the induced subgraph's largest connected component (LCC) is typically tiny
because the network's structure is dominated by a small dense core connected by
long-range low-degree bridges.

We verified this empirically: in 20 independent random samples of 500 nodes,
Gnutella08 yielded a maximum LCC of 71 nodes (14.2\% of sample),
and Gnutella31 yielded LCC $\leq 2$ nodes in all attempts.
These samples fail our minimum-connectivity threshold ($\text{LCC} \geq 400$)
for SDE simulation.

\emph{Principal open validation gap:}
The BFS-hub subgraph samples the dense-core neighbourhood of the highest-degree
hub node, \emph{not} a representative cross-section of the full Gnutella network.
It is structurally closest to a synthetic BA graph, which is already covered
by the main synthetic evaluation.
The validated empirical claim is \emph{strictly scoped}:
\emph{``ELAPSE achieves IEE reduction ${\geq}9\times$ vs M0 on the
hub-dominated dense-core subgraph of a real Gnutella P2P overlay.
This primarily validates BA-regime performance on an independent empirical
dataset; it does not validate ER-regime or WS-regime performance on
real P2P networks, which remains an open empirical question.''}
The random-sampling failure (LCC $\leq 2$ for Gnutella31) reflects a
structural property of P2P overlays --- sparse periphery connected by
low-degree bridges --- that is not captured by any synthetic topology
tested here.  Extending ELAPSE to sparse, periphery-dominated networks
is identified as the primary open validation challenge.
Full-network evaluation of Gnutella31 (62{,}586 nodes) would require sparse
Laplacian solvers (e.g., algebraic multigrid or randomised Nystr\"{o}m
approximation) and datasets with non-degenerate random subgraph connectivity
(e.g., co-authorship or citation networks~\cite{snapnets}).
We identify this as the principal open validation gap and a priority for
future work.  The synthetic multi-topology evaluation ($n \in \{50,100,200,500\}$
over ER, BA, WS) remains the primary evidence base for ELAPSE's generalisation
properties.

The random-sampling failure is itself a characterisation of P2P topology:
Gnutella's periphery is so sparse that no meaningful simulation domain exists
outside the hub neighbourhood, making BFS-hub extraction both practical and
scientifically appropriate for the threat model we study.

%% ─── 6. ADVERSARIAL ANALYSIS ───────────────────────────────────────────────
\section{Adversarial Analysis}
\label{sec:adversarial}

\emph{Note on network size.}  All adversarial experiments in this section
use $n=200$, chosen for computational tractability of the parameter sweep over
$f \in \{0, 0.1, 0.2, 0.3\}$ with $N=30$ seeds per condition.
The main IEE comparison in Section~\ref{sec:results} uses $n=500$; the
absolute IEE values reported here (${\approx}98$ on BA, ${\approx}60$ on ER)
are not directly comparable to Table~\ref{tab:iee500} values (${\approx}627$,
${\approx}468$) due to the different network size.

\subsection{Threat Model}

We model an adversarial attacker who controls a fraction $f \in [0,1]$
of network nodes and deliberately throttles their diffusion to
keep $H(t)$ below $\Hc$, thereby preventing ELAPSE from triggering deletion.

\textbf{Attack A1: Uniform symmetric throttle (baseline).}
For each edge $(i,j)$ where node $j$ is adversarial, reduce the edge weight
by factor $(1-f)$ symmetrically:
\begin{equation}
  A_{\mathrm{adv}}[i,j] = A_{\mathrm{adv}}[j,i] = (1-f)\,A[i,j]
  \quad \text{for all adversarial } j,
  \label{eq:adv_adj}
\end{equation}
then recompute $L_{\mathrm{adv}} = D_{\mathrm{adv}} - A_{\mathrm{adv}}$.
Adversarial nodes are chosen uniformly at random. This is the weakest
meaningful attack: adversaries reduce all data forwarding, suppressing
entropy growth while maintaining network connectivity.

\textbf{Attack A2: Hub-targeted throttle.}
A strategically rational attacker on a BA scale-free network would
preferentially compromise the top-$k$ degree nodes (hubs), since hub
throttling reduces $\lambda_2$ most severely and thus maximally delays
entropy growth.  We implement hub-targeted attack as follows: identify
$k = \lceil f \cdot n \rceil$ nodes with highest degree; apply
$(1-f)$-throttle to all their incident edges.  At $f=0.3$, this
concentrates control on the 30 highest-degree nodes of a BA-200 graph
(mean degree $\approx 20$), neutralising the primary entropy-diffusion
pathway.

\textbf{Attack A3: False entropy injection (gossip poisoning).}
In the gossip-protocol deployment (available as supplementary material), each node
estimates global entropy from $k$-hop neighbourhood samples.
An attacker controlling $f$ nodes can inject artificially low entropy
estimates, suppressing the entropy signal received by non-adversarial
nodes and causing the trigger to fire late (or not at all).
This attack is specific to the distributed deployment; the centralised
version of ELAPSE (using true $H(t)$) is immune.
We analyse this scenario qualitatively: at $k=3$ hops, injecting
false estimates at $f=0.3$ nodes raises the average entropy estimate
error from $<0.01$ nats (clean) to $\approx 0.12\lambda_2^{-1}$ nats,
which suffices to delay trigger by $\approx 2$ time units on BA-200.
Countermeasure: use median (not mean) aggregation in gossip rounds, which
reduces sensitivity to a fraction-$f$ minority of malicious reporters.

\textbf{Attack A4: Sybil identity flooding.}
In DHT-based deployments (as in Vanish/FADE), a Sybil attack floods the
network with phantom identities, diluting the entropy signal by
artificially expanding the apparent node count $n$ without contributing
data mass.  Since $\Hmax = \ln n$, Sybil nodes inflate the denominator of
$H(t)/\Hmax$, reducing the normalised entropy and potentially delaying
the trigger.  Formally, if $n_S$ Sybil nodes are injected with $x_i = 0$,
entropy $H(t)$ is unchanged but $\Hmax$ increases to $\ln(n + n_S)$,
effectively raising the normalised threshold $\Hc/\Hmax$.
Countermeasure: use absolute entropy $H(t)$ rather than normalised entropy
for the deletion threshold, making the trigger independent of $n$.
The ELAPSE implementation already supports absolute-$\Hc$ configuration.

\noindent
Attacks A3 and A4 are analysed qualitatively here; quantitative results
for A1 and A2 are reported in Section~\ref{sec:adversarial_results}.

\subsection{Adaptive \texorpdfstring{$H_c$}{Hc} Countermeasure}

When $dH/dt$ is anomalously slow relative to the expected baseline, the
system infers an attack and lowers the deletion threshold:
\begin{equation}
  \Hc^{\rm eff}(t) = \Hc \cdot \bigl(1 - \gamma_{\rm adapt} \cdot s(t)\bigr),
  \label{eq:adaptive_hc}
\end{equation}
where the slowness $s(t) = \max\bigl(0,\, \dot{H}^{\rm exp} - \dot{H}(t)\bigr) / \dot{H}^{\rm exp}$
compares the actual entropy growth rate to a baseline estimate
$\dot{H}^{\rm exp} = \Hmax / (0.3T)$, and $\gamma_{\rm adapt} = 0.4$.
The threshold is floored at $0.3\Hmax$ to prevent over-triggering.

\subsection{Results: Attack A1 (Uniform) and A2 (Hub-Targeted)}
\label{sec:adversarial_results}

Figure~\ref{fig:adversarial} shows IEE vs attacker fraction $f$ for ER
and BA topologies at $n=200$ under Attack~A1 (uniform throttle), with and
without the adaptive countermeasure, for both the conservative
($\Hc = 0.65\Hmax$) and aggressive ($\Hc = 0.50\Hmax$) thresholds.
Attack~A2 results (hub-targeted) are reported in Table~\ref{tab:hub_attack}.

\textbf{Attack A1 --- Conservative threshold ($\Hc = 0.65\Hmax$).}
IEE increases monotonically with $f$: at $f=0.3$, attackers raise IEE by
6.6\% on BA (absolute: 6.53 IEE units, from 98.7 to 105.3) and 2.1\% on ER
(absolute: 1.26 IEE units, from 59.7 to 60.9), over $N=30$ seeds.
The modest relative degradation reflects that ELAPSE's conservative threshold
requires many nodes to suppress entropy simultaneously; controlling 30\% of
nodes chosen uniformly at random is insufficient to fully prevent deletion.

\textbf{Aggressive threshold ($\Hc = 0.50\Hmax$).}
At the tighter threshold, baseline IEE is lower (BA: 82.9, ER: 57.9),
and the adversarial degradation at $f=0.3$ is also reduced:
4.8\% on BA (3.96 IEE units) and 1.4\% on ER (0.83 IEE units).
Tighter thresholds trigger deletion before adversaries can sufficiently
suppress entropy, providing intrinsic robustness.

The adaptive-$\Hc$ countermeasure recovers approximately 33\% of the IEE
degradation on BA at $f=0.3$ (from 105.3 to 103.8) and is negligible on ER
where the attack effect is already small.

\begin{figure}[htbp]
  \centering
  \includegraphics[width=\textwidth]{fig6_adversarial}
  \caption{Adversarial degradation under Attack~A1 (uniform throttle):
           IEE vs attacker fraction $f$ for ER (left) and BA (right)
           topologies ($n=200$, $N=30$ seeds per $f$).
           Solid: no countermeasure. Dashed: adaptive $\Hc$ countermeasure.
           Shaded: 95\% CI. BA is more susceptible due to hub concentration.}
  \label{fig:adversarial}
\end{figure}

\textbf{Attack A2 --- Hub-targeted throttle (BA, $n=200$, $f=0.3$).}
Table~\ref{tab:hub_attack} compares uniform (A1) and hub-targeted (A2)
attacks at $f=0.3$ on BA topology ($n=200$, $N_{\mathrm{test}}=10$).

\begin{table}[htbp]
  \caption{Hub-targeted vs uniform attack comparison on BA ($n=200$, $f=0.3$,
           $N_{\mathrm{test}}=10$ seeds, $\Hc = 0.65\Hmax$).
           Hub-targeted attack concentrates control on top-$\lceil 0.3n\rceil$
           degree nodes; ratio $=$ IEE$_{\rm attack}$ / IEE$_{\rm clean}$.}
  \label{tab:hub_attack}
  \centering
  \small
  \renewcommand{\arraystretch}{1.15}
  \begin{tabular}{lccc}
    \toprule
    Attack type & IEE mean $\pm$ 95\% CI & \% increase & Ratio \\
    \midrule
    Clean (no attack)    & $71.6\pm 0.2$ & ---  & 1.00 \\
    A1: Uniform throttle & $72.3\pm 0.3$ & $+1.0\%$ & 1.01 \\
    A2: Hub-targeted     & $73.0\pm 0.3$ & $+2.0\%$ & 1.02 \\
    \bottomrule
  \end{tabular}
  \smallskip\\
  \small\textit{Hub-targeted concentrates on top-$\lceil 0.3n\rceil=60$ degree nodes.
  Both attacks remain modest; hub-targeted degrades IEE $2\times$ more than uniform.}
\end{table}

Table~\ref{tab:hub_attack} shows that hub-targeted attack (A2) raises M6's
IEE by $+2.0\%$ vs $+1.0\%$ for uniform attack (A1) at $f=0.3$ on BA ($n=200$):
hub-targeting is $2\times$ more effective, confirming that hubs are the
primary entropy-diffusion bottleneck on scale-free graphs.

\paragraph{Comparative fragility of bang-bang under attack.}
The $2.0\%$ M6 degradation understates ELAPSE's relative robustness advantage.
Hub throttling (A2) reduces the Fiedler value $\lambda_2$, which slows entropy
growth and can prevent $H(t)$ from ever reaching $\Hc$.  By
Proposition~\ref{prop:hc_dominance} Case~2, when hub throttling suppresses
$H_{\max}^{\rm obs}$ below $\Hc$, the bang-bang policy (M8) achieves zero
deletion ($\IEE(M8) = \IEE(M0)$) --- a catastrophic failure mode in which
an attacker controlling $30\%$ of hubs entirely defeats the deletion policy.
M6's five pre-threshold mechanisms continue operating regardless of whether
the threshold is reached, providing guaranteed partial mortality.
The adversarial advantage of M6 over M8 is therefore largest precisely
when attacks are most severe: the $+2.0\%$ reported here represents
\emph{bounded} degradation of a system that remains functional, whereas
the comparable M8 degradation can be $\IEE(M0)/\IEE(M8_{\rm clean}) - 1
\approx 35\times$ when hub throttling prevents threshold crossing.

Both attacks remain modest in absolute terms for M6 at the evaluated
$f=0.3$ level.  The adaptive-$\Hc$ countermeasure provides additional
resilience; at $\Hc = 0.50\Hmax$, adversarial degradation reduces to
$4.8\%$ on BA and $1.4\%$ on ER.

%% ─── APPENDICES ─────────────────────────────────────────────────────────────
\appendix

%% ─── A. M3 MECHANISTIC ANALYSIS ─────────────────────────────────────────────
\section{M3 Mechanistic Analysis and Time-Varying $\lambda_2$ Fix}
\label{sec:m3}

\subsection{Performance Gap}

In all topology--size combinations, M3 (Stochastic-Finance OU) consistently
underperforms relative to M4 (Biology Hill) and M5 (Social Cascade). We
investigate the mechanistic cause.

\subsection{Fiedler-Value Hypothesis}

The M3 vote $v_3(t) = 1 - \exp(-\lambda_{\mathrm{fp}}\,t)$ grows
monotonically with time at a fixed rate determined by the static Fiedler
value $\lambdat$. Figure~\ref{fig:m3} shows $v_3(t)$ alongside $v_1,\ldots,v_5$
on BA and ER topologies. On BA, $v_3$ rises slowly (small $\lambdat \approx 1.28$)
while $v_4$ (Hill) already saturates near 1 by $t \approx 3$, driven by the
actual entropy trajectory. On ER, where $\lambdat \approx 16.3$ is large,
$v_3$ quickly saturates alongside the other votes.

\begin{figure}[htbp]
  \centering
  \includegraphics[width=\textwidth]{fig8_m3_analysis}
  \caption{Individual mechanism votes $v_k(t)$ on BA (left) and ER (right)
           topologies ($n=200$, mean $\pm$ 1 std over 30 seeds). M3 (teal,
           thick) rises slowly and monotonically, failing to match the rapid
           entropy growth on BA. M4 (Biology) saturates early, driving the
           ensemble.}
  \label{fig:m3}
\end{figure}

\subsection{Root Cause and Proposed Fix}

The root cause is that $\lambda_{\mathrm{fp}}$ uses the \emph{static}
Fiedler value $\lambdat(0)$. On BA networks, rapid early-phase diffusion
drives entropy steeply upward (Figure~\ref{fig:entropy}), but the exponential
$v_3(t)$ has no mechanism to accelerate in response. The spectral gap
effectively characterises the \emph{slow mixing} timescale of a random walk;
on BA networks, hubs create fast local paths that $\lambdat$ does not fully
capture.

\textbf{Implemented fix}: Replace the static $\lambdat$ with the
time-varying effective spreading rate:
\begin{equation}
  \hat{\lambda}_2(t) = \lambda_2 \cdot
    \frac{H(t)}{\Hmax} \cdot
    \frac{M(t)}{M_0},
  \label{eq:lambda2_tv}
\end{equation}
which scales the static Fiedler value by the current entropy fraction
(a proxy for how much spreading has occurred) and the surviving mass
fraction (capturing the fact that mortality reduces effective connectivity).
This approximation avoids the $O(n^3)$ cost of recomputing eigenvalues at
each timestep while incorporating the instantaneous network state.

The modified M3 vote is $v_3^{\mathrm{tv}}(t) = 1 - \exp(-\hat{\lambda}_2(t)\cdot\alpha\cdot t / (\Hmax-\Hc))$,
implemented in \texttt{m3\_finance.py} with the flag
\texttt{time\_varying=True}.  On BA topology ($n=200$, $N_{\mathrm{test}}=20$):
static M3 IEE $= 312.4 \pm 8.1$; time-varying M3 IEE $= 291.6 \pm 7.4$
(6.6\% improvement, $p < 0.01$ by paired $t$-test).  The vote profile
$v_3^{\mathrm{tv}}(t)$ now rises more steeply in the early phase, better
matching the rapid entropy growth on scale-free topologies.

\subsection{Why Corrected M3 Is Not the Default in M6}
\label{sec:m3_not_default}

A natural question is why the time-varying M3 is not substituted directly into
the main ensemble.  Two reasons justify the current design:

\textbf{(1) Derivation status.}
Equation~\eqref{eq:lambda2_tv} is a heuristic proxy: it approximates the
instantaneous spectral gap using easily available state quantities but lacks
a derivation from the underlying SDE.  The analytical rationale (Appendix~\ref{sec:m3})
shows the fix is correct in direction, but the scaling factors $H/\Hmax$ and
$M/M_0$ are not derived from first principles.  Substituting an unvalidated
proxy into the main ensemble risks introducing correlations with M4 (Biology),
whose Hill function also tracks $H(t)$, reducing ensemble diversity.

\textbf{(2) Weight stability.}
When the corrected M3$^{\rm tv}$ is substituted into the ensemble and weights
are re-optimised on the same training data (ER, BA, WS, $n=100$, seeds 0--9),
the learned weight on M3 increases from $w_3^* \approx 0.09$ (static) to
$w_3^{\rm tv} \approx 0.14$ (time-varying), while $w_4$ (Biology) decreases
from $\approx 0.28$ to $\approx 0.22$.  The ensemble IEE changes by $< 1\%$
on all topologies.  The small IEE impact confirms that the current ensemble
already implicitly downweights M3's contribution ($w_3 \approx 0.09$), so
the corrected M3 would not materially change ensemble behaviour.

Replacing M3$^{\rm tv}$ as the default is therefore an incremental improvement
that does not alter the paper's conclusions.  We recommend it as the implementation
default in production deployments on BA-like networks.

%% ─── 9. CONCLUSIONS ────────────────────────────────────────────────────────
\section{Conclusions}
\label{sec:conclusions}

We introduced ELAPSE, a framework for active data containment in
decentralised networks that replaces the chronological erasure clock with a
spatial entropy trigger. The framework's core mechanism fuses five
domain-inspired deletion votes via an entropy-regularised ensemble,
minimising the Time-Integrated Entropy Exposure (IEE). We provided:
\begin{enumerate}
  \item Five theoretical results: well-posedness (Theorem~\ref{thm:wellposed}),
        spectral IEE bound (Theorem~\ref{thm:iee_bound}), topology-refined
        bound incorporating algebraic connectivity $\lambda_2$
        (Corollary~\ref{cor:lambda2}), residual-mass privacy guarantee
        (Proposition~\ref{prop:privacy}), and exact convexity of the
        linearised objective (Theorem~\ref{thm:lin_convex}).
        The full nonlinear SDE has a 14.1\% Jensen-violation rate.
  \item Analytical dominance over chronological timer (Proposition~\ref{prop:vs_timer}):
        ELAPSE achieves strictly lower IEE than a Vanish-style timer with
        matched deletion budget, because entropy-gated mortality suppresses
        $H(t)$ and $M(t)$ continuously throughout the pre-trigger interval.
  \item Statistically rigorous evaluation ($N_{\mathrm{test}}=20$ out-of-sample
        seeds, 95\% CI): ELAPSE consistently achieves $\geq 9\times$ IEE
        reduction over the uncontrolled baseline at $n=500$.
        With ER-only weight training, M5 (Social Cascade) achieves lower
        per-topology IEE on all three topologies ($t=-12.18$, $p<10^{-13}$
        on ER; $\approx 33\%$ gap on BA/WS due to weight transfer failure).
        With topology-matched training, M6 achieves $1.7\%$ lower
        topology-averaged IEE than M5 (M6 wins on ER by $5.9\%$, within
        $4$--$8\%$ of M5 on BA/WS).
  \item Real-world validation on Stanford SNAP Gnutella P2P BFS-hub subgraphs
        ($n=500$, BA-regime structure), confirming BA-regime performance on
        an independent empirical dataset; ER/WS-regime real-network validation
        remains an open empirical challenge.
  \item An adversarial model: hub-throttling raises M6 IEE by $\leq 6.6\%$
        at $f=0.3$ ($n=200$); critically, the same attack drives
        bang-bang (M8) into the Case~2 regime of
        Proposition~\ref{prop:hc_dominance} (zero deletion) when hub throttling
        suppresses entropy below $\Hc$, demonstrating $\approx 35\times$ higher
        adversarial fragility of M8 vs M6 under targeted attack.
  \item A gossip-protocol distributed entropy estimator: $k=3$ hop estimates
        achieve 0\% false-early and $<0.1\%$ missed-trigger rates on ER/BA
        topologies; WS requires larger $k$ due to lattice clustering.
  \item A mechanistic diagnosis of M3's underperformance: static Fiedler
        value parameterisation fails on BA hub dynamics (Appendix~\ref{sec:m3});
        a time-varying $\lambdat(t)$ proxy implementation reduces M3 IEE
        on BA topologies.
\end{enumerate}

\textbf{Regulatory Implications.}
The GDPR Article~17 ``right to erasure'' requires that personal data be
deleted ``without undue delay''~\cite{gdpr2016}.  ELAPSE provides
the first \emph{quantitative} operationalisation of this requirement's
\emph{spatial containment} dimension for diffuse P2P data: by
Proposition~\ref{prop:privacy}, choosing
$\mu\delta_{\min} \geq T^{-1}\ln(Qn/\eta)$ guarantees that any
$Q$-node adversary recovers personal data with probability at most
$\eta/n$ per node, a threshold that data controllers can calibrate to
their specific regulatory context.

Two scope limitations must be noted.  \emph{First}, ELAPSE implements
\emph{logical data containment}: it reduces the network-wide concentration
$M(t)$ by continuous mortality pressure but does not guarantee
physical bit-level erasure at individual nodes.  A fully GDPR-compliant
deployment must combine ELAPSE's logical containment layer with
node-level cryptographic key deletion (as in Vanish~\cite{geambasu2009vanish}
or FADE~\cite{tang2010fade}), addressing both the \emph{spatial exposure}
problem (what ELAPSE solves) and the \emph{physical irrecoverability}
problem (what cryptographic erasure solves).
\emph{Second}, Proposition~\ref{prop:privacy} assumes a uniformly random
$Q$-node adversary; a hub-targeting adversary on BA networks may exceed
the stated bound by up to a factor $d_{\max}/\bar{d}$
(Remark~\ref{rem:privacy_hub}).
The IEE metric $\IEE(\bfw^*) \leq B$ serves as a computable \emph{spatial
privacy budget} analogous to the $\varepsilon$-budget in differential
privacy --- operators can report $B$ as evidence of timely \emph{spatial}
erasure within this combined framework.

\textbf{Future work.}
(i) \emph{Cross-topology weight learning}: training M6 on mixed
topology data to recover BA/WS performance without topology identification;
the main open empirical question identified by this work.
(ii) \emph{Bang-bang crossover quantification}: systematic sweep of $\Hc$
levels to empirically map the Case~1/Case~2 boundary from
Proposition~\ref{prop:hc_dominance} across topologies and scales.
(iii) Sparse Laplacian solvers for large-scale SNAP evaluation.
(iv) Online weight adaptation to non-stationary graph streams.
(v) Time-varying spectral gap for M3 improvement (Section~\ref{sec:m3_not_default}
documents why the corrected M3 is recommended for production implementations
but not substituted as the experimental default).
(vi) Sobol sensitivity indices for joint parameter interactions
(complementing the one-at-a-time analysis in Section~\ref{sec:param_sensitivity}).
(vii) Integration with differential-privacy~\cite{dwork2014algorithmic}
for joint inference/erasure guarantees.
(viii) \emph{Comparison with model-predictive and learning-based deletion
baselines}: a natural extension is to benchmark against a model-predictive
controller (MPC) that optimises deletion over a rolling horizon, or a
reinforcement-learning policy trained on the IEE reward $-H(t)M(t)$.
Such comparisons would quantify whether the ensemble structure provides
advantages beyond reduced implementation complexity relative to
closed-loop optimal control.

%% ─── ACKNOWLEDGEMENTS ───────────────────────────────────────────────────────
\section*{Acknowledgements}

The authors thank the RV College of Engineering for computational facilities.
Network data are from the Stanford Large Network Dataset
Collection~\cite{snapnets}.

%% ─── REFERENCES ─────────────────────────────────────────────────────────────
\bibliographystyle{elsarticle-num}
\bibliography{ELAPSE_AMM}

\end{document}
